% Options for packages loaded elsewhere
% Options for packages loaded elsewhere
\PassOptionsToPackage{unicode}{hyperref}
\PassOptionsToPackage{hyphens}{url}
%
\documentclass[
  letterpaper,
  8pt,
  a4paper]{book}
\usepackage{xcolor}
\usepackage{amsmath,amssymb}
\setcounter{secnumdepth}{5}
\usepackage{iftex}
\ifPDFTeX
  \usepackage[T1]{fontenc}
  \usepackage[utf8]{inputenc}
  \usepackage{textcomp} % provide euro and other symbols
\else % if luatex or xetex
  \usepackage{unicode-math} % this also loads fontspec
  \defaultfontfeatures{Scale=MatchLowercase}
  \defaultfontfeatures[\rmfamily]{Ligatures=TeX,Scale=1}
\fi
\usepackage{lmodern}
\ifPDFTeX\else
  % xetex/luatex font selection
\fi
% Use upquote if available, for straight quotes in verbatim environments
\IfFileExists{upquote.sty}{\usepackage{upquote}}{}
\IfFileExists{microtype.sty}{% use microtype if available
  \usepackage[]{microtype}
  \UseMicrotypeSet[protrusion]{basicmath} % disable protrusion for tt fonts
}{}
\makeatletter
\@ifundefined{KOMAClassName}{% if non-KOMA class
  \IfFileExists{parskip.sty}{%
    \usepackage{parskip}
  }{% else
    \setlength{\parindent}{0pt}
    \setlength{\parskip}{6pt plus 2pt minus 1pt}}
}{% if KOMA class
  \KOMAoptions{parskip=half}}
\makeatother
% Make \paragraph and \subparagraph free-standing
\makeatletter
\ifx\paragraph\undefined\else
  \let\oldparagraph\paragraph
  \renewcommand{\paragraph}{
    \@ifstar
      \xxxParagraphStar
      \xxxParagraphNoStar
  }
  \newcommand{\xxxParagraphStar}[1]{\oldparagraph*{#1}\mbox{}}
  \newcommand{\xxxParagraphNoStar}[1]{\oldparagraph{#1}\mbox{}}
\fi
\ifx\subparagraph\undefined\else
  \let\oldsubparagraph\subparagraph
  \renewcommand{\subparagraph}{
    \@ifstar
      \xxxSubParagraphStar
      \xxxSubParagraphNoStar
  }
  \newcommand{\xxxSubParagraphStar}[1]{\oldsubparagraph*{#1}\mbox{}}
  \newcommand{\xxxSubParagraphNoStar}[1]{\oldsubparagraph{#1}\mbox{}}
\fi
\makeatother


\usepackage{longtable,booktabs,array}
\usepackage{calc} % for calculating minipage widths
% Correct order of tables after \paragraph or \subparagraph
\usepackage{etoolbox}
\makeatletter
\patchcmd\longtable{\par}{\if@noskipsec\mbox{}\fi\par}{}{}
\makeatother
% Allow footnotes in longtable head/foot
\IfFileExists{footnotehyper.sty}{\usepackage{footnotehyper}}{\usepackage{footnote}}
\makesavenoteenv{longtable}
\usepackage{graphicx}
\makeatletter
\newsavebox\pandoc@box
\newcommand*\pandocbounded[1]{% scales image to fit in text height/width
  \sbox\pandoc@box{#1}%
  \Gscale@div\@tempa{\textheight}{\dimexpr\ht\pandoc@box+\dp\pandoc@box\relax}%
  \Gscale@div\@tempb{\linewidth}{\wd\pandoc@box}%
  \ifdim\@tempb\p@<\@tempa\p@\let\@tempa\@tempb\fi% select the smaller of both
  \ifdim\@tempa\p@<\p@\scalebox{\@tempa}{\usebox\pandoc@box}%
  \else\usebox{\pandoc@box}%
  \fi%
}
% Set default figure placement to htbp
\def\fps@figure{htbp}
\makeatother





\setlength{\emergencystretch}{3em} % prevent overfull lines

\providecommand{\tightlist}{%
  \setlength{\itemsep}{0pt}\setlength{\parskip}{0pt}}



 


\makeatletter
\@ifpackageloaded{bookmark}{}{\usepackage{bookmark}}
\makeatother
\makeatletter
\@ifpackageloaded{caption}{}{\usepackage{caption}}
\AtBeginDocument{%
\ifdefined\contentsname
  \renewcommand*\contentsname{Table of contents}
\else
  \newcommand\contentsname{Table of contents}
\fi
\ifdefined\listfigurename
  \renewcommand*\listfigurename{List of Figures}
\else
  \newcommand\listfigurename{List of Figures}
\fi
\ifdefined\listtablename
  \renewcommand*\listtablename{List of Tables}
\else
  \newcommand\listtablename{List of Tables}
\fi
\ifdefined\figurename
  \renewcommand*\figurename{Figure}
\else
  \newcommand\figurename{Figure}
\fi
\ifdefined\tablename
  \renewcommand*\tablename{Table}
\else
  \newcommand\tablename{Table}
\fi
}
\@ifpackageloaded{float}{}{\usepackage{float}}
\floatstyle{ruled}
\@ifundefined{c@chapter}{\newfloat{codelisting}{h}{lop}}{\newfloat{codelisting}{h}{lop}[chapter]}
\floatname{codelisting}{Listing}
\newcommand*\listoflistings{\listof{codelisting}{List of Listings}}
\makeatother
\makeatletter
\makeatother
\makeatletter
\@ifpackageloaded{caption}{}{\usepackage{caption}}
\@ifpackageloaded{subcaption}{}{\usepackage{subcaption}}
\makeatother
\usepackage{bookmark}
\IfFileExists{xurl.sty}{\usepackage{xurl}}{} % add URL line breaks if available
\urlstyle{same}
\hypersetup{
  pdftitle={BIO144 Course Information (version 2026)},
  pdfauthor={Owen Petchey},
  hidelinks,
  pdfcreator={LaTeX via pandoc}}


\title{BIO144 Course Information (version 2026)}
\author{Owen Petchey}
\date{2026-01-23}
\begin{document}
\frontmatter
\maketitle

\renewcommand*\contentsname{Table of contents}
{
\setcounter{tocdepth}{1}
\tableofcontents
}

\mainmatter
\bookmarksetup{startatroot}

\chapter*{Preface}\label{preface}
\addcontentsline{toc}{chapter}{Preface}

\markboth{Preface}{Preface}

This website contains information about the BIO144 Data Analysis in
Biology course at the University of Zurich.

Beware that Owen sometimes makes updates to the information during the
semester, so if you have downloaded a copy or taken screenshots, your
copy may not exactly match the most current version. However, any
changes will be correcting typos or improving explanations. If content
does change in an important way Owen will announce this in the lectures,
or on OLAT, or by email.

\section*{How this website was made}\label{how-this-website-was-made}
\addcontentsline{toc}{section}{How this website was made}

\markright{How this website was made}

The book was written using a type of RMarkdown. It allows a script with
a mix of normal text and R code to produce chapters and a book that has
a mixture of text, R code, and R output. Rmarkdown is very useful for
making reports, books, presentations, and even websites.

This book is a Quarto book. To learn more about Quarto books visit
\url{https://quarto.org/docs/books}.

\bookmarksetup{startatroot}

\chapter*{General information}\label{general-information}
\addcontentsline{toc}{chapter}{General information}

\markboth{General information}{General information}

\section*{Aims and learning outcomes}\label{aims-and-learning-outcomes}
\addcontentsline{toc}{section}{Aims and learning outcomes}

\markright{Aims and learning outcomes}

This course will help you develop a solid foundation in answering
biological questions with quantitative data. Moreover, the approaches
you will learn are generally applicable to using data to solve problems,
an increasing important skill in a world more and more dominated by
data.

\href{https://github.com/opetchey/BIO144/raw/master/1_learning_objectives/Weekly_learning_objectives.pdf}{The
Learning Objectives of the course} tell you what you need to learn and
know. Hence, it is very important for you to be very clear about the
learning objectives of the course. They will help you know what you're
expected to learn and what you will be examined on. If you can achieve
all of the learning objectives, you will be very well prepared for a
great future in data analysis (and for the exam).

\section*{Course schedule}\label{course-schedule}
\addcontentsline{toc}{section}{Course schedule}

\markright{Course schedule}

The course runs during the 14-week spring semester. The course has 12
weeks of new content, and 2 weeks of self-study (no new content). Each
of the 12 weeks of new content has a set of learning objectives. Each
week has four parts: a lecture, homework, a practical session, and a
practice quiz.

Lectures are on Monday mornings, 8:00 - 9:45. Practical sessions are on
Thursday or Friday afternoons, 13:00 - 14:45.

You'll be assigned to a practical session on either Thursday or Friday.
If at any point you wish to change afternoons permanently, contact
\texttt{studienkoordination@biol.uzh.ch}. If you want to swap days for
just one or two weeks, please do so without asking.

Please see the ``Year specific information'' section below for a
detailed course schedule.

\section*{What should I already know?}\label{what-should-i-already-know}
\addcontentsline{toc}{section}{What should I already know?}

\markright{What should I already know?}

Please see the section in OLAT ``Previous Knowledge''.

\section*{Expected workload}\label{expected-workload}
\addcontentsline{toc}{section}{Expected workload}

\markright{Expected workload}

The course is worth 4 credit points, so with 1 CP being approximately 30
hours work, requires around 120 hours study. Approximately one third of
this is ``in class'' and two thirds is self-study, including in-course
assessments and preparation for the examination.

\section*{Bring your own laptop}\label{bring-your-own-laptop}
\addcontentsline{toc}{section}{Bring your own laptop}

\markright{Bring your own laptop}

Please bring your laptop to class. There are some power outlets in the
lecture theatre, but best will be to make sure that your laptop has
sufficient battery charge to last the duration of the class (up to 2
hours), and consider using an external additional battery (powerbar?) if
required. Please make sure you know how to use your laptop (some people
are sometimes not so sure, perhaps because its new, or they borrow
someone else's.) Your laptop must be able to connect to the wireless
internet at UZH.

We will be using three computer programmes during class (as well as the
usual suspects, such as an internet browser etc.) please make sure you
have them:

\begin{enumerate}
\def\labelenumi{\arabic{enumi}.}
\item
  A spreadsheet programme, such as Microsoft Excel, Mac Numbers, or
  OpenOffice.
\item
  \href{https://www.r-project.org}{The R software environment for
  statistical computing and graphics}. It's worth knowing why we are
  using R. You can find out by reading Section 4 of the Preface of the
  Beckerman, Childs, and Petchey Book.
\item
  \href{https://posit.co/downloads/}{Rstudio}. We use RStudio to
  interact with R. Rstudio has lots of features that make working with R
  much much easier and enjoyable.
\end{enumerate}

If you need any help setting up your laptop with R and RStudio, please
look at relevant section of the ``Getting R \& RStudio'' section of the
``Previous Knowledge'' section of the course. If you still have
problems, ask during the first practical.

\section*{What to do each week}\label{what-to-do-each-week}
\addcontentsline{toc}{section}{What to do each week}

\markright{What to do each week}

\begin{itemize}
\item
  Attend the lecture.
\item
  After the lecture and before the practical do the homework.
\item
  Come to and do the practical.
\item
  Do the practice quiz.
\item
  Have fun ;)
\end{itemize}

\section*{Course texts}\label{course-texts}
\addcontentsline{toc}{section}{Course texts}

\markright{Course texts}

The required content of the course is all available in the course book:
\href{https://opetchey.github.io/BIO144_Course_Book/}{BIO144 Course
Book}. The lectures will help you understand the content of the book.
The homework and practicals will also help you learn the material.
Practice quizzes will help you check your learning, and prepare for the
final exam. All of this will also help you develop your skills in R,
which will also be required in the final exam.

There is other optional reading that we may suggest to you during the
course. This is listed below. Textbooks for optional reading are in the
main library, and should be accessible electronically through the main
library.

Optional reading comes from the following books. Also each chapter of
the course book suggest reading to further your understanding.

\begin{itemize}
\tightlist
\item
  \href{https://github.com/opetchey/BIO144/raw/master/9_assets/Stahel_LinReg.pdf}{\emph{Linear
  Regression, Seminar für Statistik, ETH Zürich}(2008 / 2013) Werner
  Stahel.} And here is an
  \href{https://github.com/opetchey/BIO144/raw/master/9_assets/Stahel_LinReg_EN.pdf}{English
  version of that text}.
\item
  \href{https://academic.oup.com/book/27290}{\emph{Getting Started with
  R, An introduction for biologists} (2017, Second Edition) Beckerman,
  Childs \& Petchey.} (\textbf{DO NOT USE THE FIRST EDITION!}). The link
  is to the book on the Oxford University Press website; full access to
  the book will require you have a UZH ip address (i.e.~are on the UZH
  network or VPN into it).
\item
  \href{https://academic.oup.com/book/39565}{Insights from Data} with R,
  by Petchey, et al.~(Same access information as the previous book in
  this list).
\item
  \emph{The New Statistics with R, An Introduction for Biologists}
  (2015, First Edition) Hector.
  \href{https://academic.oup.com/book/32474}{Ebook via UZH library.}
\end{itemize}

Some other nice (though optional) texts that you might find interesting.

\begin{itemize}
\item
  \emph{Regression; Models, Methods and Applications} (2013 Edition)
  Fahrmeir, Kneib, Lang \& Marx.
\item
  \emph{The Analysis of Biological Data} (2015, Second Edition) Whitlock
  \& Schluter.
\item
  The Essential Guide to Effect Sizes. Statistical Power, Meta-Analysis,
  and the Interpretation of Research Results (2010, First Edition)
  Ellis.
  \href{http://ezproxy.uzh.ch/login?url=http://dx.doi.org/10.1017/CBO9780511761676}{Ebook
  via UZH library.}
\item
  \emph{Statistics: An Introduction using R} (2011, Second Edition)
  Crawley.
\end{itemize}

\section*{Getting datasets}\label{getting-datasets}
\addcontentsline{toc}{section}{Getting datasets}

\markright{Getting datasets}

You can get the datasets used in the course
\href{https://github.com/opetchey/BIO144/tree/master/3_datasets}{here}.
To download one dataset from that web page, click on the dataset and you
then see a webpage with that dataset. Then right click on the ``Raw''
button, and select ``Save As'', select location to save to, and press
Save. Alternatively you can get a
\href{https://github.com/opetchey/BIO144/raw/master/3_datasets/all_datasets_bio144.zip}{zip
file of all the datasets here}, which you will then need to unpack
before use.

You can get datasets associated with the book Getting Started with R, by
Beckerman, Childs, and Petchey
\href{http://r4all.org/books/datasets/}{here}.

IMPORTANT: if you open the downloaded file in Excel, perhaps to have a
look at it, then \textbf{\emph{DO NOT, UNDER ANY CIRCUMSTANCES, ASK OR
ALLOW EXCEL TO SAVE IT}}. (Excel does not like R and will do everything
in its power to mess things up for R. Don't let it do that.)

Here is the data you collected during the first lecture:

\href{https://lms.uzh.ch/url/RepositoryEntry/17654940136/CourseFolder/0/path\%3D~~human\%5Fbenchmark\%5Fdata\%5FFS25\%2Ecsv/0}{human\_benchmark\_data\_fs25.csv}

\section*{Really important!!!}\label{really-important}
\addcontentsline{toc}{section}{Really important!!!}

\markright{Really important!!!}

\textbf{Please please please read, understand, and try your very best to
remember what is below.}

We will very often be importing data into R using the read.csv or
read\_csv function. In the time since the material of the course was
written (including texts such as Getting Started with R) the makers of R
changed something important about how importing data works. The best
solution is for us to make sure we change it back to the old way, so all
the course material makes sense!

Here are two things you must always do:

When using read.csv always add the argument `as.is = FALSE.\\
For example: dd \textless- read.csv(``my data file.csv'', as.is = FALSE)

When using read\_csv, always pipe the result into
mutate\_if(is.character, factor).\\
For example:\\
dd \textless- read\_csv(``my data file.csv'') \%\textgreater\%\\
mutate\_if(is.character, factor)

\section*{Attendance}\label{attendance}
\addcontentsline{toc}{section}{Attendance}

\markright{Attendance}

\textbf{Attending class is a really good idea!}: ``attendance {[}is{]} a
better predictor of college grades than any other known predictor of
academic performance'', Credé, Roch, \& Kieszczynka, 2010.
\href{https://journals.sagepub.com/doi/abs/10.3102/0034654310362998}{The
full article.}

We give you the freedom to use the learning materials and opportunities
in the way that suits you best. We strongly believe that attending the
practical classes and lectures will be an efficient, fun, and valuable
way for you to learn.

We will not be tracking attendance. If you do not come to class
(practical or lecture) the only penalty you will potentially experience
is loss of some learning opportunity.

\textbf{What if I was ill or otherwise unable to attend?} Please use
your initiative to catchup. Lectures will be available as podcasts. Ask
for help from friends for help with practicals. And whatever else you
can think of.

\section*{Self study weeks}\label{self-study-weeks}
\addcontentsline{toc}{section}{Self study weeks}

\markright{Self study weeks}

\subsection*{Consolidate, review, reinforce,
expand!}\label{consolidate-review-reinforce-expand}
\addcontentsline{toc}{subsection}{Consolidate, review, reinforce,
expand!}

During approximately two weeks there is no lecture and no practical.
During these week you should take the opportunity to consolidate,
review, reinforce, expand!

Here is some optional reading you could look at during those weeks:

\begin{itemize}
\tightlist
\item
  \href{Amrhein_etal2019.pdf}{Retire Statistical Significance}
\item
  \href{significanceVSimportance_Whitlock.pdf}{Interleaf from Whitlock
  and Schluter: Statistical significanc vs.~biological importance}
\item
  \href{correlationVScausation.pdf}{Interleaf from Whitlock and
  Schluter: Correlation vs.~causation}
\item
  \href{NZZ_April2016.pdf}{P-Werte in der NZZ (3. April 2016)} (You will
  have to zoom in and probably read it on the screen!)
\item
  And finally, a really thoughtful article in one of the world's leading
  scientific journals: \href{Goodman_2016_Science.pdf}{S. Goodman
  (2016): Aligning statistical and scientific reasoning}
\end{itemize}

\section*{Assessment}\label{assessment}
\addcontentsline{toc}{section}{Assessment}

\markright{Assessment}

There will be two assessment-related types of activity:

\subsection*{1. Weekly practice quizzes}\label{weekly-practice-quizzes}
\addcontentsline{toc}{subsection}{1. Weekly practice quizzes}

These are only for you. Your scores do not contribute to anything. The
are intended:

\begin{itemize}
\tightlist
\item
  To make sure you know if you're learning and understanding the
  material.
\item
  To motivate you to keep up with the content of the class, and thereby
  reduce the chance of you falling behind and then having trouble in the
  final exam.
\item
  To give you practice with the types of questions that will be in the
  final examination.
\end{itemize}

There is no deadline for the practice quizzes. You can do them whenever
you like. You can do them multiple times if you like.

\subsection*{2. The final exam (and the resit
exam)}\label{the-final-exam-and-the-resit-exam}
\addcontentsline{toc}{subsection}{2. The final exam (and the resit
exam)}

A final assessment during the module examination period. Your
performance in this is the only determinant of your final grade.

\textbf{Examination date and time:} Please see the Studium Biology,
Modulprufungen im Grundstudium web page and the links therein for the
examination dates and other important examination information.

\textbf{\emph{This information is draft for 2026 and may still be
changed.}}

Further information:

\begin{itemize}
\tightlist
\item
  The exam will contain about 40 questions. Each will be one of three
  types: single choice, multiple choice, or numeric. Single choice is a
  question with several possible answer, of which only one is correct --
  you indicate which. Multiple choice is a question with several
  possible answers; none, 1, 2, 3, or 4 may be correct -- you indicate
  which. Numeric requires you to enter a number to answer the question.
\item
  A full mock exam will be available on OLAT.
\item
  During the exam you are allowed \textbf{three sheets of A4 paper} that
  you may put information on both sides of. You are also permitted to
  bring a dictionary (it must have no writing in it). No other resources
  are permitted.
\item
  You are allowed to bring into the exam one translation dictionary
  (Übersetzungswörterbuch) (e.g., an Englisch--Deutsch-Wörterbuch /
  English--German dictionary). Examination staff should be allowed to
  inspect this resource.
\item
  You will need to present your UZH card during the exam, to verify your
  identity.
\item
  You will need to use your laptop continuously during the two hours of
  the exam. Some power outlets are available, but best is to make sure
  you have enough battery power to last the exam. Please ensure your
  laptop can connect to the internet via wifi.
\item
  You will take the exam in a special instance of OLAT from within the
  \emph{Safe Exam Browser} that must be installed on your computer.
\item
  From within the Safe Exam Browser you will be able to use an RStudio
  Server.
\item
  There will be a mandatory practice exam experience including use of
  the Safe Exam Browser. Details will be sent to you.
\item
  The exam is in English.
\item
  You must agree with an Honour Code on OLAT immediately before you take
  the exam.
\item
  You will sit the exam in a lecture theatre on the Irchel Campus.
\item
  The final exam requires you to analyse some datasets that are provided
  to you. Some students have reported that they didn't have enough time
  during the exam, so please ensure that you can do things like
  importing data quickly and easily.
\item
  Your working (e.g.~any R code your write) is not assessed directly. It
  is assessed indirectly by whether you answer a question correctly.
\item
  During the exam, if you have any trouble importing the data, please
  immediately ask for help.
\item
  When answering numeric questions, the default is to give your answer
  with precision of 2 decimal places. Some questions specify to give the
  answer to a different number of decimal places. Please be careful to
  take care of this.
\end{itemize}

\section*{Important information about the ``practice exam
experience''}\label{important-information-about-the-practice-exam-experience}
\addcontentsline{toc}{section}{Important information about the
``practice exam experience''}

\markright{Important information about the ``practice exam experience''}

The practice examination experience is only to give you experience of
the examination interface, the types of questions that can be asked, and
to test the examination interface. Do not use the number and type of
different types of question as an indicator of anything about the final
exam.

\href{https://github.com/opetchey/BIO144/raw/master/2_content/xtras/BIO144_exam_preparation_solutions.pdf}{Here
are the questions and solutions that are used in the practice exam
experience}.

You will be informed by email when the practice exam experience is, and
all further details.

\section*{Getting help}\label{getting-help}
\addcontentsline{toc}{section}{Getting help}

\markright{Getting help}

We aim that you can ask for help about anything and get a useful answer
as soon as possible. The following guidelines will help you get fast,
useful help. The primary ways to ask for help are:

\begin{enumerate}
\def\labelenumi{\arabic{enumi}.}
\tightlist
\item
  During the practicals on Thursday and Friday. You will get close to
  instant responses.
\item
  Anytime on the Discussion Forum of this OLAT course. We aim that you
  get a response in less than two days. Anonymous posts are possible.
\item
  After lectures.
\end{enumerate}

Of course, please feel free to email Owen directly if you wish. Please
don't hesitate to ask for help!

\subsection*{Posting code problems on the Discussion
Forum}\label{posting-code-problems-on-the-discussion-forum}
\addcontentsline{toc}{subsection}{Posting code problems on the
Discussion Forum}

The Wiki:FAQ contains a page \textbf{\emph{How to ask a question so we
can easily help}}. Please take a look.

\subsection*{Specialist statistical
assistance}\label{specialist-statistical-assistance}
\addcontentsline{toc}{subsection}{Specialist statistical assistance}

I (Owen) am not a statistician. Hence, I might be unable to answer some
of your questions about more statistical issues, particularly as we
experience more and more complex statistics. To deal with this,
Prof.~Furrer has agreed to help with about statistics that I cannot
answer (or cannot answer fully). Please post them as a new thread in the
discussion forum. Include as much information as possible, including R
code and output as appropriate. Please do not contact Prof.~Furrer
directly. I will then contact Prof.~Furrer and feedback the answer to
you.

\section*{Giving feedback}\label{giving-feedback}
\addcontentsline{toc}{section}{Giving feedback}

\markright{Giving feedback}

Your views on the course are really important to us. What works well,
what doesn't work so well, what content is particularly important, what
content is not so important, what content is missing?

There are several ways of providing feedback.

\begin{enumerate}
\def\labelenumi{\arabic{enumi}.}
\tightlist
\item
  During the course, you can tell us in person during class, and can
  make posts on the Forum. Such feedback during the course has the
  potential to improve the course for you. I.e., we might be able to
  change something quickly, so your and your fellow students' experience
  of the class is improved.
\item
  Fill in the ``official'' course evaluation form. Details to follow.
\item
  If have anonymous feedback to give us before or after the official
  course evaluation then please write on some paper and post it through
  the University internal mail to us.
\end{enumerate}

\section*{Course language}\label{course-language}
\addcontentsline{toc}{section}{Course language}

\markright{Course language}

The course will be mostly in English. Many of the TAs speak and
understand English and German. Prof.~Petchey best understands and speaks
English.

Please let us know how we can assist with any challenges this presents.

\href{http://isi.cbs.nl/glossary/index.htm}{Here is a web page that
gives most common (and many not so common) statistical terms in a
multitude of languages!}

\section*{Reference letters}\label{reference-letters}
\addcontentsline{toc}{section}{Reference letters}

\markright{Reference letters}

The course instructors and the TAs try to get to know each of you during
the course, particularly in the practical classes. That said, it is
unlikely we gain enough experience of you to write for you a meaningful
letter of reference / recommendation. Better will be to develop
relationships during subsequent block courses that can lead to
meaningful letters.

\section*{Example solution scripts}\label{example-solution-scripts}
\addcontentsline{toc}{section}{Example solution scripts}

\markright{Example solution scripts}

Example solution scripts can be found here at the link below. However,
please control your use of these example scripts carefully. Students in
previous years have said they felt that they sometimes looked at the
solutions to early or too quickly. They are provided because many
students over many years have asked for them.

\url{https://github.com/opetchey/BIO144/tree/master/6_example_scripts}

\section*{Lecture podcasts}\label{lecture-podcasts}
\addcontentsline{toc}{section}{Lecture podcasts}

\markright{Lecture podcasts}

All lectures will be recorded and made available for later viewing.

\bookmarksetup{startatroot}

\chapter*{Year specific information}\label{year-specific-information}
\addcontentsline{toc}{chapter}{Year specific information}

\markboth{Year specific information}{Year specific information}

\section*{2026 Detailed schedule}\label{detailed-schedule}
\addcontentsline{toc}{section}{2026 Detailed schedule}

\markright{2026 Detailed schedule}

Lecture: Monday, 8:00-9:45am) Practical: Thurs or Fri. 1-3pm)

\begin{longtable}[]{@{}ll@{}}
\toprule\noalign{}
Week beginning & Content \\
\midrule\noalign{}
\endhead
\bottomrule\noalign{}
\endlastfoot
Feb.~16 & U1 Introduction, review, demonstration \\
Feb.~23 & U2 R \& RStudio \\
Mar.~2 & U3 Regression part 1 \\
Mar.~9 & U4 Regression part 2 \\
Mar.~16 & U5 Analysis of variance \\
Mar.~23 & U6 Multiple regression \\
Mar.~30 & U7 Interactions \\
Apr.~6 & -- No lecture, no practicals \\
Apr.~13 & U8 Count data \\
Apr.~20 & -- No lecture, review / catch-up practicals \\
Apr.~27 & U9 Binary data \\
May 4 & U10 Ordination \\
May 11 & U11 Mixed models \& What next? \\
May 18 & U12 Review \\
May 25 & -- No lecture, review / catch-up practicals \\
\end{longtable}

\section*{2026 Locations}\label{locations}
\addcontentsline{toc}{section}{2026 Locations}

\markright{2026 Locations}

Lectures are in Y24-G-45\\
(Lectures will be live-streamed and podcasts will be made available, in
any case.)

Practicals are in Y24-G-45.

\section*{2026 Assessment}\label{assessment-1}
\addcontentsline{toc}{section}{2026 Assessment}

\markright{2026 Assessment}

Practice examination experience: Date and time to be communicated

Final Examination date and time: Please see this Studium Biology,
Modulprufungen im Grundstudium web page and the links therein for the
examination dates and other important examination information.

\section*{2026 Practice exam}\label{practice-exam}
\addcontentsline{toc}{section}{2026 Practice exam}

\markright{2026 Practice exam}

We will share a practice exam with you to experience the exam platform
and types of questions in the final exam. We likely make it available
and notify you after the spring break.

\section*{2026 Exam viewing}\label{exam-viewing}
\addcontentsline{toc}{section}{2026 Exam viewing}

\markright{2026 Exam viewing}

If you did not get a 4.0 or higher:

You can view your final examination paper with Dr Frank Pennekamp -
date, time and location to be announced by email.

You can view your resit examination paper with Dr Frank Pennekamp -
date, time and location to be announced by email.

\bookmarksetup{startatroot}

\chapter*{Learning Objectives}\label{learning-objectives}
\addcontentsline{toc}{chapter}{Learning Objectives}

\markboth{Learning Objectives}{Learning Objectives}

This document describes the learning objectives for \textbf{BIO144: Data
Analysis in Biology}.

\begin{center}\rule{0.5\linewidth}{0.5pt}\end{center}

\section*{Course-level learning
objectives}\label{course-level-learning-objectives}
\addcontentsline{toc}{section}{Course-level learning objectives}

\markright{Course-level learning objectives}

After successfully completing BIO144, students will be able to:

\begin{enumerate}
\def\labelenumi{\arabic{enumi}.}
\tightlist
\item
  Formulate biological questions as statistical models, identifying
  appropriate response and explanatory variables.
\item
  Use R and RStudio to import, explore, visualise, and analyse
  biological data reproducibly.
\item
  Fit, interpret, and compare statistical models commonly used in
  biology, including linear models and generalised linear models.
\item
  Evaluate model assumptions, diagnose violations, and understand the
  consequences for inference.
\item
  Interpret model output quantitatively and biologically, rather than
  mechanically reporting p-values.
\item
  Communicate statistical results clearly, using figures, tables, and
  written explanations appropriate for biological audiences.
\item
  Read and critically assess quantitative analyses in published
  biological research.
\end{enumerate}

\begin{center}\rule{0.5\linewidth}{0.5pt}\end{center}

\section*{Chapter-level learning
objectives}\label{chapter-level-learning-objectives}
\addcontentsline{toc}{section}{Chapter-level learning objectives}

\markright{Chapter-level learning objectives}

\subsection*{Introduction}\label{introduction}
\addcontentsline{toc}{subsection}{Introduction}

After this chapter, students will be able to:

\begin{itemize}
\tightlist
\item
  Explain the goals, structure, and expectations of the BIO144 course.
\item
  Know the prior knowledge and skills required for success in the
  course.
\item
  Understand appropriate use of AI tools in data analysis work and
  during the course.
\item
  Describe a typical data analysis workflow, from developing question to
  communicating the answer.
\end{itemize}

\begin{center}\rule{0.5\linewidth}{0.5pt}\end{center}

\subsection*{R and RStudio}\label{r-and-rstudio}
\addcontentsline{toc}{subsection}{R and RStudio}

After this chapter, students will be able to:

\begin{itemize}
\tightlist
\item
  Navigate the RStudio interface and explain the purpose of scripts, the
  console, and the environment.
\item
  Run R code and understand basic R syntax.
\item
  Import data into R and inspect its structure.
\item
  Know the many ways one can get help with R and RStudio.
\item
  Understand what are and how to work with add-on packages in R.
\item
  Perform basic data manipulation operations, including importing,
  viewing, and various manipulations.
\item
  Produce simple exploratory data visualisations.
\item
  Use scripts to support reproducible data analysis workflows.
\end{itemize}

\begin{center}\rule{0.5\linewidth}{0.5pt}\end{center}

\subsection*{Simple linear regression -- Part
1}\label{simple-linear-regression-part-1}
\addcontentsline{toc}{subsection}{Simple linear regression -- Part 1}

After this chapter, students will be able to:

\begin{itemize}
\tightlist
\item
  Explain the purpose of simple linear regression in biological data
  analysis.
\item
  Understand the mathematical form of a simple linear regression model.
\item
  Describe how the slope and intercept are calculated.
\item
  Describe how the error / variation is modelled.
\item
  State the assumptions underlying simple linear regression.
\item
  Assess if the assumptions are reasonably met using diagnostic plots.
\item
  Recognise and remedy common problems encountered when fitting linear
  regression models.
\item
  Perform linear regression in R.
\end{itemize}

\begin{center}\rule{0.5\linewidth}{0.5pt}\end{center}

\subsection*{Simple linear regression -- Part
2}\label{simple-linear-regression-part-2}
\addcontentsline{toc}{subsection}{Simple linear regression -- Part 2}

After this chapter, students will be able to:

\begin{itemize}
\tightlist
\item
  Understand how to measure how good is the regression (correlation and
  R-squared).
\item
  Test and explain if the parameter estimates are compatible with some
  specific value (t-test).
\item
  Understand how to find the range of parameters values are compatible
  with the data (confidence intervals).
\item
  Understand what are the regression lines compatible with the data
  (confidence band), and how to construct these in R.
\item
  Understand what are the plausible values of newly collected data
  (prediction band) and how to calculate these in R.
\item
  Interpret the biological meaning of regression coefficients.
\item
  Know how to communicate the results of a linear regression analysis
  effectively.
\item
  Critically assess the strength and limitations of linear regression
  models.
\end{itemize}

\begin{center}\rule{0.5\linewidth}{0.5pt}\end{center}

\subsection*{One-way ANOVA}\label{one-way-anova}
\addcontentsline{toc}{subsection}{One-way ANOVA}

After this chapter, students will be able to:

\begin{itemize}
\tightlist
\item
  Explain how ANOVA fits within the linear model framework.
\item
  Fit a one-way ANOVA model using \texttt{lm()}.
\item
  Interpret group means and differences between groups.
\item
  Explain the concept of variance partitioning.
\item
  Understand the connection between ANOVA and regression with
  categorical predictors.
\item
  Decide when ANOVA is an appropriate modelling approach for biological
  data.
\item
  Perform linear regression in R, assess model assumptions, and
  interpret results in a biological context.
\item
  Communicate ANOVA results effectively using text, tables, and
  visualisations.
\end{itemize}

\begin{center}\rule{0.5\linewidth}{0.5pt}\end{center}

\subsection*{Multiple regression}\label{multiple-regression}
\addcontentsline{toc}{subsection}{Multiple regression}

After this chapter, students will be able to:

\begin{itemize}
\tightlist
\item
  Understand what is a question that multiple regression can help
  answer.
\item
  Know the mathematical form of multiple regression models.
\item
  Fit linear models with multiple explanatory variables.
\item
  Assess if model assumptions are reasonably met using diagnostic plots.
\item
  Assess if the group of explanatory variables significantly explain
  variation in the response.
\item
  Describe which variables are associated with the response, and the
  direction of these associations.
\item
  Measure the amount of variation explained by the model (R-squared,
  adjusted R-squared).
\item
  Assess the relative importance of different explanatory variables.
\item
  Make conditional predictions from multiple regression models.
\item
  Understand what is collinearity, and how it affects multiple
  regression.
\item
  Use R to fit, diagnose, and interpret multiple regression models.
\item
  Communicate multiple regression results effectively using text,
  tables, and visualisations.
\end{itemize}

\begin{center}\rule{0.5\linewidth}{0.5pt}\end{center}

\subsection*{Interactions}\label{interactions}
\addcontentsline{toc}{subsection}{Interactions}

After this chapter, students will be able to:

\begin{itemize}
\tightlist
\item
  Explain what an interaction between explanatory variables means
  biologically and mathematically.
\item
  Fit models that include interaction terms.
\item
  Interpret interaction coefficients correctly.
\item
  Visualise and explain how effects depend on the values of other
  variables.
\item
  Distinguish additive from non-additive biological effects.
\item
  Judge when interactions are biologically meaningful and justified.
\item
  Understand interactions in the context of both categorical and
  continuous explanatory variables, ANCOVA, two-way ANOVA, and multiple
  regression.
\item
  Communicate results involving interactions clearly and accurately.
\item
  Use R to fit, diagnose, interpret, and communicate models with
  interactions.
\end{itemize}

\begin{center}\rule{0.5\linewidth}{0.5pt}\end{center}

\subsection*{Generalised linear models for count
data}\label{generalised-linear-models-for-count-data}
\addcontentsline{toc}{subsection}{Generalised linear models for count
data}

After this chapter, students will be able to:

\begin{itemize}
\tightlist
\item
  Describe the kinds of biological questions and processes that involve
  count data.
\item
  Recognise when, why, and how count data violate the assumptions of
  linear models.
\item
  Explain why normal error assumptions are inappropriate for count data.
\item
  Describe the key components of a generalised linear model.
\item
  Fit Poisson regression models using \texttt{glm()}.
\item
  Interpret model coefficients on the appropriate scale.
\item
  Assess whether a Poisson model provides an adequate description of the
  data.
\item
  Use R to fit, diagnose, and interpret GLMs for count data.
\item
  Communicate results from count data analyses effectively.
\end{itemize}

\begin{center}\rule{0.5\linewidth}{0.5pt}\end{center}

\subsection*{Generalised linear models for binary
data}\label{generalised-linear-models-for-binary-data}
\addcontentsline{toc}{subsection}{Generalised linear models for binary
data}

After this chapter, students will be able to:

\begin{itemize}
\tightlist
\item
  Identify binary response variables in biological datasets, and be
  familiar with common examples (e.g., presence--absence,
  success--failure) and the types of questions they can arise from.
\item
  Explain why linear models are unsuitable for binary outcomes.
\item
  Fit logistic regression models using \texttt{glm()}.
\item
  Interpret model coefficients in terms of probabilities and odds.
\item
  Visualise fitted relationships for binary data.
\item
  Understand how logistic regression supports biological inference about
  presence--absence or success--failure outcomes.
\item
  Use R to fit, diagnose, and interpret GLMs for binary data.
\item
  Communicate results from binary data analyses effectively.
\end{itemize}

\begin{center}\rule{0.5\linewidth}{0.5pt}\end{center}

\subsection*{Ordination and multivariate
data}\label{ordination-and-multivariate-data}
\addcontentsline{toc}{subsection}{Ordination and multivariate data}

After this chapter, students will be able to:

\begin{itemize}
\tightlist
\item
  Recognise when biological questions involve many response variables
  simultaneously.
\item
  Explain the goals of ordination methods.
\item
  Interpret ordination plots in terms of similarities and differences
  among observations.
\item
  Understand ordination as a tool for dimensionality reduction.
\item
  Critically assess what information ordination methods do and do not
  provide.
\item
  Use R to perform basic ordination analyses (i.e., PCA, NMDS).
\item
  Communicate ordination results effectively using visualisations and
  text.
\end{itemize}

\begin{center}\rule{0.5\linewidth}{0.5pt}\end{center}

\section*{Mixed models and what next}\label{mixed-models-and-what-next}
\addcontentsline{toc}{section}{Mixed models and what next}

\markright{Mixed models and what next}

\subsection*{Part 1: Mixed models}\label{part-1-mixed-models}
\addcontentsline{toc}{subsection}{Part 1: Mixed models}

After completing this part, students will be able to:

\begin{itemize}
\tightlist
\item
  Recognise biological situations in which data are \textbf{grouped,
  nested, or hierarchical}.
\item
  Explain why standard linear models may be inappropriate for grouped
  data.
\item
  Describe the conceptual difference between \textbf{fixed effects} and
  \textbf{random effects}.
\item
  Understand random effects as a way of modelling \textbf{structured
  sources of variation}.
\item
  Interpret mixed models as extensions of linear models, rather than as
  fundamentally different tools.
\item
  Identify common biological examples where mixed models are appropriate
  (e.g.~repeated measures, individuals within populations, sites within
  regions).
\item
  Interpret mixed model output at a conceptual level, focusing on
  biological meaning rather than technical detail.
\item
  Appreciate the role of \textbf{partial pooling} in balancing
  information across groups.
\item
  Recognise the limitations and assumptions of mixed models without
  needing to master implementation details.
\item
  Use R to fit simple linear mixed models using the \texttt{lme4}
  package.
\item
  Communicate results from mixed model analyses effectively, focusing on
  biological interpretation.
\end{itemize}

\begin{center}\rule{0.5\linewidth}{0.5pt}\end{center}

\subsection*{Part 2: What next?}\label{part-2-what-next}
\addcontentsline{toc}{subsection}{Part 2: What next?}

After completing this part, students will be able to:

\begin{itemize}
\tightlist
\item
  Place the statistical models learned in BIO144 within a broader
  landscape of data analysis methods.
\item
  Recognise when more advanced models may be required to address
  biological questions.
\item
  Identify directions for further learning in statistics and data
  analysis.
\item
  Understand that statistical modelling is an \textbf{iterative and
  evolving process}, not a fixed set of rules.
\item
  Reflect critically on the limits of the models used in the course.
\item
  Appreciate the importance of biological reasoning alongside
  statistical tools.
\item
  Approach future quantitative methods courses and analyses with
  confidence and curiosity.
\end{itemize}

\begin{center}\rule{0.5\linewidth}{0.5pt}\end{center}

\section*{Review (L12)}\label{review-l12}
\addcontentsline{toc}{section}{Review (L12)}

\markright{Review (L12)}

No additional learning objectives.

\bookmarksetup{startatroot}

\chapter*{FAQs}\label{faqs}
\addcontentsline{toc}{chapter}{FAQs}

\markboth{FAQs}{FAQs}

\section*{Organization-related FAQ}\label{organization-related-faq}
\addcontentsline{toc}{section}{Organization-related FAQ}

\markright{Organization-related FAQ}

\subsection*{How to ask a question so we can easily
help?}\label{how-to-ask-a-question-so-we-can-easily-help}
\addcontentsline{toc}{subsection}{How to ask a question so we can easily
help?}

You're probably quite competent at asking questions so that we can
easily help. But here are a few tips that can help you be even better:

\begin{itemize}
\item
  If you're asking a question in the OLAT Forum and include R code,
  please try to ensure it is easy to read and copy into R.
\item
  If you're asking for help about some code that doesn't work, please
  consider making a ``minimal, reproducible example''. This web page
  describes what one is and how to make one
  (\url{https://stackoverflow.com/help/minimal-reproducible-example}).
\end{itemize}

\subsection*{Is practical attendance
mandatory?}\label{is-practical-attendance-mandatory}
\addcontentsline{toc}{subsection}{Is practical attendance mandatory?}

You are free to do the practicals whenever you want, but we encourage
that you do attend the practical afternoons, especially if you have
questions!

\subsection*{Can I change which practical afternoon I
attend?}\label{can-i-change-which-practical-afternoon-i-attend}
\addcontentsline{toc}{subsection}{Can I change which practical afternoon
I attend?}

In general it's better that you stick to the afternoon you registered
for, but if you cannot attend it you can of course join the practical on
the other day.

\subsection*{Why is the number of attempts for the questions
limited?}\label{why-is-the-number-of-attempts-for-the-questions-limited}
\addcontentsline{toc}{subsection}{Why is the number of attempts for the
questions limited?}

For now, the tests are implemented such that you can take each test
\textbf{only} once, but you can suspend every test you take. This saves
your progress within the test and when you get back to it you'll pick up
exactly where you left it (your progress is saved). You can only take
the tests once, because if you can take it more than once your progress
is lost, meaning that you would have to retake even the questions that
you already answered correctly (OLAT does not allow us to change this).

However, within each test attempt, you get \textbf{2 tries per
question}:

\begin{itemize}
\tightlist
\item
  For the weekly graded assessments there must be a limited number of
  tries for the graded assessment questions.
\item
  For self-tests in the homework and practical parts this is not set in
  stone. For now, for almost all questions you get 2 tries because
  giving more attempts might demotivate some people. This might change
  in the future.
\end{itemize}

\subsection*{Are the sections named Extras part of the
exam?}\label{are-the-sections-named-extras-part-of-the-exam}
\addcontentsline{toc}{subsection}{Are the sections named Extras part of
the exam?}

No, they are just for your interest!

\subsection*{Device with enough battery for
exam}\label{device-with-enough-battery-for-exam}
\addcontentsline{toc}{subsection}{Device with enough battery for exam}

If you are not certain that the battery of your laptop can last for the
duration of the final examination, please let us know using the form at
the end of the course content here on OLAT.

\subsection*{About the Discussion
Forum}\label{about-the-discussion-forum}
\addcontentsline{toc}{subsection}{About the Discussion Forum}

Use the discussion forum to ask questions, give tips, make suggestions
about how the course could be improved\ldots{} or anything else that
springs to mind while you're taking the course.

However, before posting please check out the Index in the ``Wiki: FAQ''
section in this course to see whether your question has already been
answered there and also look at the other forum posts here to see
whether there is already a discussion thread for your topic. If you
cannot find what you are looking for, go ahead and post something in
this forum, we will respond as soon as possible!

Other potential uses of the forum are that you may make a list of
questions that you want to ask during the class, or a list of things
that you'd like us to go over again in class, or in more detail in
lectures or practicals. Then we can look at these in lectures or
practicals.

One rule: do not post answers to assessment or solutions to exercises.
Help others find and realise for themselves the answer, but don't post
the actual answer. If you post comments/questions about graded
assessments, we will always investigate and take appropriate action.
This may include deleting or editing your post. Please ask about GAs
during a practical or email.

Another rule / guideline. Open a new post / thread for a new issue.

\section*{R-related FAQ}\label{r-related-faq}
\addcontentsline{toc}{section}{R-related FAQ}

\markright{R-related FAQ}

\subsection*{What version of R is
needed?}\label{what-version-of-r-is-needed}
\addcontentsline{toc}{subsection}{What version of R is needed?}

It is good practice to keep R up to date, so it is recommended that you
have the most recent version of R installed. In any case, you should at
least have the R version 4.0 or newer installed. You can check the
version by running ``version'' or ``sessionInfo()'' in R.

\subsection*{What packages are needed?}\label{what-packages-are-needed}
\addcontentsline{toc}{subsection}{What packages are needed?}

Note: this list is incomplete and will be updated in the future.

The packages you need in each unit are usually in the lecture and/or in
the homework material.

An incomplete list (work in progress):

\begin{itemize}
\tightlist
\item
  tidyverse (contains ggplot2, readr, dplyr, tidyr, etc.)
\item
  skimr
\item
  ggfortify
\item
  \ldots{}
\end{itemize}

\subsection*{I am having issues installing a
package}\label{i-am-having-issues-installing-a-package}
\addcontentsline{toc}{subsection}{I am having issues installing a
package}

Problems with installing packages are somewhat common. Often, these
problems be avoided by using

\texttt{install.packages(\ "NAME",\ type\ =\ "binary")}

instead of simply

\texttt{install.packages(\ "NAME",\ )}

The reason is that quite often, packages are available a source
packages, which must be compiled locally, which is often not possible on
standard R machines. Therefore, by specifying type = ``binary'', you
tell R to install the already compiled version, even though it might be
an older version.

\subsection*{Different output of summary()/str()/etc. function than what
is shown in
video}\label{different-output-of-summarystretc.-function-than-what-is-shown-in-video}
\addcontentsline{toc}{subsection}{Different output of
summary()/str()/etc. function than what is shown in video}

During the course it will happen that the output that you get from
function such as summary(), glimpse(), str(), \ldots, will look a bit
different from the output that you see in the (homework) videos.

Usually, the reason for this is that since R version 4.0.0 (and newer)
the function read.csv() no longer reads in non-numerical variables as
factor variables, but instead these variables are read in as character
variables and functions such as summary(), glimpse(), str(), \ldots{}
handle character variables differently than they handle factor
variables.

\subsection*{What does rm(list=ls()) do?}\label{what-does-rmlistls-do}
\addcontentsline{toc}{subsection}{What does rm(list=ls()) do?}

While working in R we assign stuff to variables, load in data, etc.
These things are stored in the R ``environment''. The line rm(list=ls())
deletes everything that is stored in this environment. The reason for
this is basically that by freeing up space we are on top of things (when
many different things are stored in the environment it can get messy)
and are sure that we are not using things from (for instance) a
previously run script.

The function rm() stand for ``remove'' so it removes everything that we
give it. In this case we give it a list of things: ls() lists every name
that is stored in the environment. We usually use it as the first thing
in a script.

\subsection*{How do I set the number of decimal places
shown?}\label{how-do-i-set-the-number-of-decimal-places-shown}
\addcontentsline{toc}{subsection}{How do I set the number of decimal
places shown?}

Use the round() function to tell R how many decimal places to report and
what rounding to do. You otherwise cannot know for sure how a number
displayed in the console has been rounded or not.

\subsection*{What does the predict() function
do?}\label{what-does-the-predict-function-do}
\addcontentsline{toc}{subsection}{What does the predict() function do?}

In general, the predict() function takes as argument a model object, for
instance one fitted with lm(). It then predicts/estimates/calculates the
response variable based on the values of the covariates (explanatory
variables) and the estimated coefficients according to the fitted model.
You can do this by hand (good for understanding).

If you do not specify the data to use for the prediction, it will
automatically use the data that was used in fitting the model. You can
also specify the dataset used, which is useful, for instance when you
have multiple covariates and you want to visualize the relation of one
specific explanatory variable with the response variable. In that case
you create a new dataset in which you keep every covariate at constant
values (e.g.~their mean) with the exception for the one in you are
interested in. Then you can use this dataset in the predict() function
and then, if you want, you can plot the relation of that variable with
the response.

Note: predict() does \textbf{not} fit a model, it evaluates a model with
some given data

\subsection*{What does the expand.grid() function
do?}\label{what-does-the-expand.grid-function-do}
\addcontentsline{toc}{subsection}{What does the expand.grid() function
do?}

The expand.grid() function is a handy function to create new data, as it
returns every possible combinations of the input provided. For instance
expand.grid(a = 1:2, b = 4:5) returns

\begin{verbatim}
    a b
1   1 4
2   2 4
3   1 5
4   2 5
\end{verbatim}

We often use this function to create new data because we want to
visualize the effect of one particular variable on the response variable
according to the fitted model. When we have a model with more than 1
explanatory variable, these variables can vary together and if they do
the effect of one particular variable on the response variable is more
difficult to see. Hence to best show the importance of a particular
variable, we keep all other variables at constant values. These constant
values should ideally be sensible values, i.e.~keeping the mass at 0
does not make sense, but it makes more sense to keep at the mean value.

\subsection*{AIC(model) versus
extractAIC(model)}\label{aicmodel-versus-extractaicmodel}
\addcontentsline{toc}{subsection}{AIC(model) versus extractAIC(model)}

You can find an explanation of the difference
\href{https://stats.stackexchange.com/questions/43733/what-is-the-difference-between-aic-and-extractaic-in-r}{here}
(see the accepted answer)

It does not matter which of the two functions you use, but be sure to
always use the same function within an analysis/script!

\subsection*{The autoplot function does not show the leverage
plot}\label{the-autoplot-function-does-not-show-the-leverage-plot}
\addcontentsline{toc}{subsection}{The autoplot function does not show
the leverage plot}

When there is one categorical and one continuous variable (and in other
cases as well), the forth plot shown by the autoplot() function (in the
ggfortify package) is not the leverage plot that we would usually get.
See also
\href{http://r4all.org/posts/diagnostic-plots-of-models-with-categotical-explanatory-variables/}{here}.
It has to do with how this function is implemented.

If you want to see the usual leverage plot, see the above link to see
how to achieve that.

The other three plots shown by the function are more important, so it is
ok that we do not see the leverage plot sometimes.

\subsection*{Cooks distance in leverage
plot}\label{cooks-distance-in-leverage-plot}
\addcontentsline{toc}{subsection}{Cooks distance in leverage plot}

I am not aware that autoplot() has a method to show Cook's distance.
This is not so important. The important thing about that plot is to see
whether there are data points that have both a large leverage (a big
impact on the regression fit result) and large residual value. These
points would be in the top and bottom right corner of the fourth plot.
If there are such points it can be interesting to repeat the analysis
without these points to see how the result changes. Note that this plot
is not needed to check the model assumptions, so for now at least the
other three plots are more important.

\subsection*{gather() vs pivot\_longer() and spread() vs
pivot\_wider()}\label{gather-vs-pivot_longer-and-spread-vs-pivot_wider}
\addcontentsline{toc}{subsection}{gather() vs pivot\_longer() and
spread() vs pivot\_wider()}

This course was written when \texttt{gather} and \texttt{spread} were
still the go-to \texttt{tidyr} functions for these tasks (there are base
\texttt{R} functions for them as well). Since then,
\texttt{pivot\textbackslash{}\_wider} and
\texttt{pivot\textbackslash{}\_longer} have been released, which have
more functionalities and are more flexible. For such simple tasks both
work perfectly fine, with the advantage that \texttt{spread} and
\texttt{gather} have slightly more intuitive names and are those more
easily understood. But of course, as recommended in the help page, the
newer functions should be used moving forward (at some point this course
will switch as well).

\subsection*{unique() vs levels()}\label{unique-vs-levels}
\addcontentsline{toc}{subsection}{unique() vs levels()}

The function \texttt{levels()} only works with factors. In newer
versions of \texttt{R} variables are read in by
\texttt{read\textbackslash{}\_csv} as characters instead of factors (see
also
\href{https://lms.uzh.ch/auth/RepositoryEntry/17168236698/CourseNode/104350296466146/path\%3DDifferent\%5Foutput\%5Fof\%5Fsummary\%28\%29~~str\%28\%29~~Etc.\%5Ffunction\%5Fthan\%5Fwhat\%5Fis\%5Fshown\%5Fin\%5Fvideo/0}{here}).
Use the function \texttt{unique()} instead of \texttt{levels()}

\section*{Statistics-related FAQ}\label{statistics-related-faq}
\addcontentsline{toc}{section}{Statistics-related FAQ}

\markright{Statistics-related FAQ}

\subsection*{What is the difference between fitted and predicted
values?}\label{what-is-the-difference-between-fitted-and-predicted-values}
\addcontentsline{toc}{subsection}{What is the difference between fitted
and predicted values?}

\textbf{Fitted values}: The response values you get when you use the
observed values of the explanatory variables, i.e.~the ones you used to
fit the model: the original data

\textbf{Predicted values}: In general, the response values you get when
you use new observations of the explanatory variables (i.e.~new data).
Note: if the new data is the same as the original data, then predicted
values == fitted values.

This being said, it happens that the terms fitted values and predicted
values are used interchangeably, even though there is a (small)
difference.

\subsection*{Interpretation of
p-values}\label{interpretation-of-p-values}
\addcontentsline{toc}{subsection}{Interpretation of p-values}

The \emph{p}-value is best interpreted as:

The \emph{p}-value is the probability of having a test statistics value
at least as extreme as the value actually observed (e.g.~the
\emph{t}-value of the \emph{t}-test or \emph{F}-value of the
\emph{F}-test) given that the Null hypothesis is correct (often H0:
parameter=0). If the p-value is below 0.05 we reject H0 and go with H1
(often H1: parameter =/= 0). The p-value is \textbf{not} the probability
of H0 being correct.

See also the lecture material.

\subsection*{Interpretation of confidence
intervals}\label{interpretation-of-confidence-intervals}
\addcontentsline{toc}{subsection}{Interpretation of confidence
intervals}

Let's assume we have a simple linear regression with body-fat as the
response and BMI as the explanatory variable. We find that the slope of
BMI is 1.75 with a corresponding 95\% confidence interval
{[}1.55,1.97{]}. This is interpreted as follows:

\begin{itemize}
\tightlist
\item
  ``All true values for beta1 (slope for BMI) between 1.55 and 1.97 are
  compatible with the observed data''.
\item
  To rephrase this slightly, this means: the \textbf{true} slope for BMI
  could have any value between 1.55 and 1.97 and we would not be
  surprised to see the data that we have (and thus, if the \textbf{true}
  slope was outside the interval, it would be surprising to see the data
  that we have: the estimated slope based on the data does not match the
  \textbf{true} slope.
\end{itemize}

\textbf{Note}: I want to stress that it is \textbf{not correct} to say
that there is a 95\% probability that the \textbf{true} slope is between
1.55 and 1.96 (see below).

So, with the interpretation given above, if the confidence interval of
the slope does not cover 0, then a slope of 0 is not compatible with the
data. In fact, in the case of the 95\% confidence interval this means
that the p-value is below 0.05, which means that we reject the Null
hypothesis H0 that the slope for BMI is 0. \textbf{Note}: this is
neither ``good'' nor ``bad'', it is just what it is: we are just
interested in the truth.

There is an alternative interpretation of confidence intervals which is
arguably more precise (but you can use the above for this course). To
explain it, here are a couple of things that I personally think are good
to keep in mind:

\begin{itemize}
\tightlist
\item
  The parameter of interest (e.g.~the slope of BMI) is estimated based
  on the available data, meaning that the estimate is a function of the
  data. If for instance the experiment is repeated the data will look a
  little bit different and the estimate will be a little bit different.
  However: the true value of the parameter of interest (the slope) is
  the same in both experiments!
\item
  The same thing goes for confidence intervals: they are estimated based
  on the data, so if the data changes a bit, the upper and the lower
  bound of the confidence interval will be a bit different as well.
\item
  An example of wrong interpretation: continuing the example above,
  let's say we know the true slope of BMI and it is 2. What is the
  probability that the confidence interval includes the \textbf{true}
  slope? 0 (!), because 2 is not in {[}1.55,1.97{]}. And what would the
  probability be if the \textbf{true} slope is 1.6? 1 (!), because 1.6
  is included in {[}1.55,1.97{]}. \textbf{This is why it is incorrect to
  say that there is a 95\% probability that the estimated confidence
  interval contains the true slope}: the \textbf{estimated} confidence
  interval either contains or it does not contain the true slope.
\end{itemize}

Now, because of probability theory reasons, a confidence interval is
best interpreted like this:

\begin{itemize}
\tightlist
\item
  ``If we repeat the entire experiment many times (and thus collect
  slightly different data each time) and for each experiment we
  calculate the 95\% confidence interval, then 95\% of the calculated
  confidence intervals will contain the true parameter value''.
\item
  This definition is a bit harder to get, but it basically tells us that
  it is a measure of uncertainty of the estimated slope.
\end{itemize}

\subsection*{Model diagnostics residuals versus fitted
values}\label{model-diagnostics-residuals-versus-fitted-values}
\addcontentsline{toc}{subsection}{Model diagnostics residuals versus
fitted values}

This entry is not complete

The central assumption of linear regression is that the errors are
independent (from each other) and identically distributed (i.i.d.)
following a normal distribution with mean = 0 and variance = sigma\^{}2.
We cannot measure the errors directly, but we have estimates of them:
the residuals. This means that we can check these assumptions based on
the residuals, because given that they are estimates of the errors and
that the model is valid the residuals have to be i.i.d. following a
normal distribution as well. In the figure residuals vs.~fitted values
we can check for some of these assumptions:

\begin{itemize}
\tightlist
\item
  Mean = 0: the mean has to be zero across all fitted values. So the
  points in the figure have to be spread approximately symmetrically
  above and below 0, so that the mean is approximately 0 across all
  fitted values. If for instance we see that for larger fitted values
  the mean of the residuals increases or decreases clearly away from 0,
  then this assumption might be violated.
\item
  Variance == sigma\^{}2: constant variance. Similarly to above, the
  variance has to be independent from the fitted values, meaning that
  the spread of point below and above 0 needs to be approximately the
  same across all fitted values. If this is not the case this results in
  a pattern (e.g.~larger fitted values result in a higher variance in
  the residuals) and this assumption might be violated.
\end{itemize}

\subsection*{Difference between binomial and count
data}\label{difference-between-binomial-and-count-data}
\addcontentsline{toc}{subsection}{Difference between binomial and count
data}

For count data we use Poisson regression and for binomial data we use
binomial regression (commonly called logistic regression)

\begin{itemize}
\tightlist
\item
  Count data come from questions like ``how many times have you voted in
  your life?'' and the answer can be a positive integer (including 0),
  while binary data come from questions like ``did you vote last year?''
  to which the answer can only be yes (1) or no (0).
\item
  We get binomial data if we have n (=number of persons) independent
  Bernoulli trials (i.e.~binary trials), so we would basically count how
  many people have voted last year out of the n we asked. In other
  words, if we have 7 yes and 3 no to the question ``did you vote last
  year?'', they come from 10 different people (sample size n=10), while
  if to the question ``how many times have you voted in your life?'' we
  get an answer like ``I've voted 10 times in my life'' this comes from
  a single person (sample size n=1), so keep this in mind.
\end{itemize}

\subsection*{Dispersion parameter}\label{dispersion-parameter}
\addcontentsline{toc}{subsection}{Dispersion parameter}

Why is the value of the estimated dispersion parameter different to
residual deviation divided by the residual degrees of freedom?

The actual dispersion is based on a weighted version of the residuals:

object is the output of a glm( \ldots{} ``pseudopoisson'') call:

dispersion \textless-
sum(object\(weights*object\)residuals\^{}2)/object\$df.residual

residuals are not:

deviance.resid \textless- sum(residuals(object, type =
``deviance'')\^{}2)

E.g.

counts \textless- c(18,17,15,20,10,20,25,13,12)

outcome \textless- gl(3,1,9)

treatment \textless- gl(3,3)

object \textless- glm(counts \textasciitilde{} outcome + treatment,
family = ``quasipoisson'')

summary(object)

sum(object\(weights*object\)residuals\^{}2)/object\$df.residual

sum(residuals(object, type = ``deviance'')\^{}2)

Next question: why weighted residuals, and how are they weighted :)

\subsection*{Attenuation in linear
regression}\label{attenuation-in-linear-regression}
\addcontentsline{toc}{subsection}{Attenuation in linear regression}

In the measurement error lecture in unit 12 there is a method for
calculating the true slope from the naive (including measurement error)
slope given on the following slide. Please see the ``Motivating
example'' section of the
\url{https://en.wikipedia.org/wiki/Errors-in-variables_models} page on
Wikipedia for a derivation of the equation show on the slide.

\href{/Image:attenuationslide.png}{\pandocbounded{\includegraphics[keepaspectratio]{attenuationslide.png}}}

\section*{Other student questions from previous years
(unsorted)}\label{other-student-questions-from-previous-years-unsorted}
\addcontentsline{toc}{section}{Other student questions from previous
years (unsorted)}

\markright{Other student questions from previous years (unsorted)}

\subsection*{Gender Imbalance Effect on
t-test}\label{gender-imbalance-effect-on-t-test}
\addcontentsline{toc}{subsection}{Gender Imbalance Effect on t-test}

\textbf{Student question (2021)}: I wanted to ask a general question
related to both the lecture and practical session: how much does gender
imbalance influence the results of a t-test? Is there a technique to
test for this possible effect and, in this case, how could one account
for the imbalance?

\textbf{Answer given}:

In general, an imbalanced dataset can be used for a t-test as long as
the assumptions of the test are met and there are enough datapoints in
each of the two compared groups (but in theory the test also works with
very small sample sizes).

The assumptions of the t-test are that the datapoints are independently
and identically distributed following a normal distribution. So as long
as there is no evidence for a violation of these assumptions, you can
use the test (you will see in the coming weeks how to check whether the
assumptions are met or not!).

By the way, usually the t-test assumes equal variances between the two
groups. If they are not (and you set var.equal=FALSE in R) the test will
have to estimate the variance in both groups and hence there will be
fewer degrees of freedom left.

Regarding how to counteract data imbalance, one possible option could be
to create a balanced dataset by randomly subsampling from the more
abundant group. But as I've said, it is not needed in this case.

More stuff to think about: it is maybe even more important to think
about whether the data is representative about the population that we
want to infer something from. I do not know what the gender ratio is in
the class, but it could be that the data imbalance is there because one
gender was more willing to participate than the other. If that is the
case, then there could be a confounding variable that influenced both
the outcome (the reaction time) and the willingness to participate. For
instance it could be that in one gender only persons with a fast
reaction time participated, while in the other group everyone
participated regardless of reaction times and thus the dataset would not
be representative of the actual class and the t-test result no longer
valid.

\textbf{Imbalanced}: means that the two groups that are compared consist
of an unequal number of datapoints

\subsection*{How to get extreme values}\label{how-to-get-extreme-values}
\addcontentsline{toc}{subsection}{How to get extreme values}

Student question (2021):

In Unit2 Question3 Exercise1 it is written that it would be complicated
to progamatically check if the extreme values belong to one single
individual. I think I have found a way to do so:

extremes \textless- slice(df,-c(1:252)) \#generate empty data frame for
(i in colnames(df)) \{ print(filter(df,df{[}i{]}==max(df{[}i{]}))) \} df
is my dataframe.

So my question is if this is a correct approach and if it would lead to
the correct answer. And my second question is, how can I add the output
of the statement to the empty data frame I created above? In python I
would solve it with a list or something similar but I am not quite sure
how I should or can do it in R.

Looking forward to an answer.

Answer given:

You found a good start to do it! Here is how you could add it to the
dataframe:

extremes \textless- df \%\textgreater\% slice(0) for (i in colnames(df))
\{ max \textless- filter(df,df{[}i{]}==max(df{[}i{]})) extremes
\textless- rbind(extremes, max) \} rbind() stands for ``row bind'' and
does just that: it binds rows together. There is also cbind(), for
columns.

Note: It is worth pointing out that in this way you ``only'' get the
maximum value of every given variable. However, It is a matter of how
extremes are defined: it could be that the max value of a variable is
not an extreme because it is perfectly in line with the other values.
Similarly, there could be more than 1 extreme value for a given
variable, with respect to the other values. The same goes for small
values, as those can be extremes as well.

More info: a common definition of extreme values is: smaller than the
first quartile minus 1.5IQR (the interquartile range) OR larger than the
third quartile plus 1.5IQR (the interquartile range). The outliers in
boxplots are for instance calculated in this way. This is not needed for
what we are doing in the course and what you did is perfectly
appropriate!

\subsection*{Lecture 3 linear regression Hypothesis and p
value}\label{lecture-3-linear-regression-hypothesis-and-p-value}
\addcontentsline{toc}{subsection}{Lecture 3 linear regression Hypothesis
and p value}

\textbf{Student question (2021):} I am a bit confused\ldots{} In lecture
3 on slide 31 it is written: „In included is the assumption that the
data follow the simple linear regression model.`` Then on slide 35 is is
written that such a little p-value suggests that it is very unlikely to
see such an slope if there would not be a correlation. But I have heard
several times before this course that a small p-value suggests to reject
the Null Hypothesis. Therefore I am confused because shouldn't the Null
Hypothesis be that there is no correlation and because of the small
p-value we accept because there is a correlation?

Can somebody help me?

\textbf{Answer given:} I think that you understood it correctly, but you
just got confused by that sentence on slide 31.

What is meant with ``In H0 included is the assumption that the data
follow the simple linear regression model'' is that in addition to H0
there is the assumption that the data can be analyzed with the chosen
regression model. It can be analyzed like this if the modelling
assumptions (see slide 25 in lecture 3) are met (you will see in the
coming weeks how to check whether they are met or not). Only if the
modelling assumptions are met it makes sense to test the actual H0.

As you correctly understood, the actual Null hypothesis H0 is that the
slope beta = 0. At the same time, if the slope is 0 it also means that
the correlation would be 0. But for completeness sake I also want to
mention that there is a difference between the regression slope and the
correlation, which is nicely explained here.

\subsection*{Plotting Regression line}\label{plotting-regression-line}
\addcontentsline{toc}{subsection}{Plotting Regression line}

\textbf{Student question (2021)}: dear BIO144 team, I tried to graph my
results using code below. However, when I tried to draw the regression
line using ``geom\_smooth'' ,it gives me 5 different regression line in
different colors for each country (could not upload the graph here).
What if I want only one regression line for all my data, and at the same
time different colors for the continents points?

\begin{verbatim}
ggplot(health\_filtered, aes(x=logExp, y=logMort, colour=continent))+
   geom\_point(size=0.8, alpha=0.9)+
   geom\_smooth(method="lm", se= FALSE)+
   xlab("Log of Child Mortality")+
   ylab("Log of healthcare Expenditure")+
   ylim(0,3)+
   xlim(1,4)+
   theme\_light()
\end{verbatim}

\textbf{Answer given}: You have specified the aesthetic mapping colour =
continent in the ggplot function. Any aesthetic mapping specified in the
ggplot function are inherited by all other geoms. I.e. they apply
automatically to all other geoms, unless we say otherwise. So your
geom\_smooth is also using the colour = continent aesthetic mapping, and
so is doing a separate regression for each.

If we want different colours for each point, but one regression, then
here are two solutions:

\begin{itemize}
\tightlist
\item
  Remove the colour = continent aesthetic mapping from the ggplot
  function, and put it only in the geom\_point function.
\end{itemize}

OR

\begin{itemize}
\tightlist
\item
  Remove the inherited colour = continent aesthetic mapping from the
  geom\_smooth by adding mapping = aes(colour = NULL). i.e.~use
  geom\_smooth(method=``lm'', mapping = aes(colour = NULL))
\end{itemize}

\subsection*{Interpreting the table lecture 4 page 29
2021}\label{interpreting-the-table-lecture-4-page-29-2021}
\addcontentsline{toc}{subsection}{Interpreting the table lecture 4 page
29 2021}

\textbf{Student question (2021)}:

The conclusions I get from interpreting the table is that

\begin{itemize}
\tightlist
\item
  bmi has a strong correlation with bodyfat (steep slope)
\item
  age probably doesn't have that strong of an influence on bodyfat (only
  0.13 slope)
\item
  the p-values all show significance but in the case of age it might not
  play an important biological role
\item
  the confidence interval of the intercept is rather wide so it could be
  useful to generate more data points (?)
\end{itemize}

Did I miss anything or did I get it wrong?

\textbf{Answer given}:

Hi, solid interpretation! Some comments:

\begin{itemize}
\tightlist
\item
  You rise a good point of statistical significance versus effect size
  (and corresponding clinical significance, in this case). It can happen
  that a coefficient is estimated to be significantly different from the
  null effect, but at a same time the estimated coefficient is so small
  (i.e.~small effect size) that it does not matter.
\item
  However, it is not always so easy to determine whether an effect size
  is big or small and in the case of the covariates \emph{age} and
  \emph{bmi} I argue that both might be important, and here is why. At
  first glance the two estimated slopes are very different, with the one
  for \emph{bmi} being roughly 14 times as big as the one for
  \emph{age}. Yet, \emph{age} is a continuous variable and the slope
  means that for every year that passes the \emph{bodyfat} percentage is
  estimated to increase by 0.13. Hence, after roughly 14 years the
  \emph{bodyfat} percentage is predicted to increase by 1.82, roughly
  the same as if the \emph{bmi} would have increased by 1. What I want
  to say is that 0.13 might not seem much, but it adds up over the
  years. On the other hand, an increase of 1 of the \emph{bmi} is
  actually not so little and thus it might be less likely to happen.
\end{itemize}

Regarding the intercept:

\begin{itemize}
\tightlist
\item
  In general more data would increase the precision of the estimates
  (i.e.~decrease the standard errors), but it is often not possible to
  have more data.
\item
  Think about what the intercept is in this model. Are we interested in
  it at all? The answer is ``probably not'': it is the \emph{bodyfat}
  percentage for when \emph{age=0} and \emph{bmi=0}, both of which are
  not possible. This is also why the estimated intercept does not make
  sense (-31 \emph{bodyfat} percentage?). In this model we are only
  interested in the estimated slopes.
\end{itemize}

In any case, good start with interpretation! You could further try to
interpret the confidence intervals (in the sense: what do they mean?).

More info: related to the point above, another thing to keep in mind
when interpreting coefficients are the units of the variables. For
\emph{age} it's years, so as I wrote above one unit increase in
\emph{age} (i.e.~1 more year) results in an expected increase of
\emph{bodyfat} percentage by 0.13. If however the unit of the variable
age was \emph{decades} instead of years, then the slope would have been
estimated to be (about) 10\emph{0.13=1.3 and in this case it would have
roughly been the same as for }bmi* and our first ``impression'' of its
size would have been different.

\subsection*{Overdispersion Index}\label{overdispersion-index}
\addcontentsline{toc}{subsection}{Overdispersion Index}

\textbf{Student question (2021)}: In the GSWR book it says only to worry
about dispersion if the dispersion index is above 2, but in the mock
exam it said that there is overdispersion even though the index was
below 2. Does that mean there is always overdispersion if the index is
above 1 but you just don't worry about it?

\textbf{Answer given}: Hi. Yes if the the dispersion parameter is
\textgreater1 it is overdispersion, but how much larger than 1 it has to
be to become problematic is up for debate. I cannot put it better than
what is written in the book:

``\emph{What if the dispersion index had been 1.2, 1.5, 2.0, or even 10?
When should you start to worry about overdispersion? That is a tricky
question to answer. One common rule of thumb is that when the index is
greater than 2, it is time to start worrying. Like all rules of thumb,
this is only meant to be used as a rough guide. In reality, the worry
threshold depends on things like sample size and the nature of the
overdispersion.}''

\subsection*{IC10 Abalone - quasipoisson and
cor}\label{ic10-abalone---quasipoisson-and-cor}
\addcontentsline{toc}{subsection}{IC10 Abalone - quasipoisson and cor}

\textbf{Student question (2021)}: I was wondering why there's no AIC
when doing the dropterm function for a full model with
family=quasipoisson? dropterm(full\_glm\_quasi, sorted=TRUE)

Also I did not quite understand the last line in the sample solution
script:

\begin{verbatim}
 cor(cbind(x=fitted(noViscera), y=fitted(lm\_model)))
\end{verbatim}

how does this work and what does the result mean?

\textbf{Answer given}: The short answer is this: you might remember that
to be calculated, the AIC needs the likelihood. I don't think that we
formally defined the likelihood in this course, so it's ok if you do not
really know what it is. What is important here is that for the GLM with
quasi-Poisson the likelihood is no longer defined and thus the AIC
cannot be calculated. I do not actually know whether \texttt{dropterm}
is useful in this case (note that the exercise does not ask you to use
it with the quasi-Poisson model, it only asks you to change the selected
Poisson model to a quasi-Poisson model).

Regarding the second question: with

\begin{verbatim}
`cor(cbind(x=fitted(noViscera), y=fitted(lm\_model)))` 
\end{verbatim}

a matrix with two columns is passed to \texttt{cor}, and from the help
page of \texttt{cor} we know that in this case the correlations between
the columns are calculated (x vs x, x vs y, y vs x and y vs y, hence 4
values). The values on the diagonal of the produced correlation matrix
are obviously 1, because we correlated respectively x vs x and y vs y.
The values on the off-diagonals are the one of interest and are the same
because x vs y and y vs x is the same. As there are only 2 columns in
this case, I find it easier to just use
\texttt{cor(fitted(noViscera),\ fitted(lm\textbackslash{}\_model))},
i.e.~to just pass the 2 vectors separately, in which case only the
correlation between the two vectors is calculated. As the correlation
between the fitted values of the two models is almost 1 (it is
\texttt{0.9733156}) you rightfully state that there is no practical
difference in the predicted values between the two models (but the glm
is still the model that should be used because the response is count
data).

\subsection*{Exercise 7c correlation between
parameters}\label{exercise-7c-correlation-between-parameters}
\addcontentsline{toc}{subsection}{Exercise 7c correlation between
parameters}

\textbf{Student question (2021)}: the scatter plot of the Betas
indicates a possible correlation between B0 and B1 as well as between B0
and B2. On the other hand, B1 and B2 do not seem to correlate. can you
help me understand why.

Is it correct to interpret this the following way: keeping one point
fixed (say B2), then the variation in the intercept will define the
variation in B1?

\textbf{Answer given}: This is what I think is going on:

\begin{itemize}
\tightlist
\item
  First, B1 and B2 are not correlated because x1 and x2 are not
  correlated either, so their coefficients are independent from each
  other.
\item
  Second, by adding random noise to the response y it can happen that
  the apparent relation between e.g.~x1 and x weakens (slope B1 closer
  to 0) or strengthens (the opposite). At the same time because the data
  (x1) did not change if the slope (B1) changes then the intercept (B0)
  changes as well, otherwise the regression line would not fit the data.
  It's not so easy to explain, maybe it helps if you try to draw it
  (i.e.~add more noise and see what happens to the intercept and the
  slope). The same goes for B2 and B0.
\end{itemize}

\subsection*{Question about Linear Model from Practical
4}\label{question-about-linear-model-from-practical-4}
\addcontentsline{toc}{subsection}{Question about Linear Model from
Practical 4}

\textbf{Student question (2021):}

I have a question about the practical we did in week 4, the one with the
milk data set.

So there we used the following linear model:

lm(formula = kcal.per.g \textasciitilde{} neocortex.perc + mass, data =
milkdat)

Now inf we check the table you gave us in the ``useful table'' section
from Lesson 7, this would constitute Case D: same slope and different
intercepts.

Now in the Practical from week 4 one question is:

What is the estimated slope of the relationship between kcal.per.g and
mass?

And the answer is -0.0054

Second question is: What is the estimated slope of the relationship
between kcal.per.g and neocortex.perc?

And the answer is 0.018

So obviously there's a contradiction here. Because in the table it says
if we use the ``+'' in the linear model, so a model without interaction,
the slopes for the two relationships between explanatory variables and
response variable should be the same. But this is not the case here.

So is there a mistake I do in thinking? I would be very glad if you
could clear this up!

\textbf{Answer given:}

The ``useful table'' is based on an ANCOVA (topic of week 7), i.e.~there
is a continuous variable (\texttt{density}) and a categorical variable
(\texttt{season}). The slope is for \texttt{density} and in the case of
no interaction (i.e.~a ``+'' between the explanatory variables) the
slope associated with \texttt{density} is the same for all values of
\texttt{season}. For the categorical variable, \texttt{k-1} (\texttt{k}
being the number of levels of the variable) intercepts are estimated.

In the milk example, \texttt{neocortex.perc} and\texttt{mass} are both
continuous variables, meaning that for both a slope is estimated
separately. Hence the questions are ``What is the estimated slope of the
relationship between kcal.per.g and \textbf{mass?}'' and ``What is the
estimated slope of the relationship between kcal.per.g and
\textbf{neocortex.perc}?''

Note that an interaction between two continuous variables would also be
possible. It would mean that the slope of one continuous variable on the
response variable changes as the values of a second continuous change.

\textbf{Follow-up question}:

So does this mean that if I have both continuous variables, if I do a
linear model with ``+'' (no interaction) in the summary table I get one
intercept for the alphabetically first explanatory variable, then the 2
slopes of the 2 explanatory variables? And in the case of ANCOVA (hence
one continuous variable and one categorical variable) I get the
intercept for the alphabetically first explanatory variable, and then
the slope which is the same, and then the difference between the first
intercept (the reference) and the intercept of the 2nd explanatory
variable (alphabetically second)? Is this correct?

\textbf{Answer given:}

Almost\ldots{} The intercept is not for the continuous variables. It
would be what you say it is if you would have a categorical variable. I
suggest that you try to write down the model to help you interpret the
summary table until you're more confident with it. In this case with
y=kcal.per.g and x1=neocortex.perc and x2=mass the model is:

y\_i =\beta\_0 + \beta\_1x\_i\^{}\{(1)\} +
\beta\_2x\_i\^{}\{(2)\}+\epsilon\_i So you are estimating 1 intercept
(beta0) and 2 slopes (one for each explanatory variable, beta1 and
beta2). Note that the intercept is basically the value of y\_i (minus
the error) when both x1=0 and x2=0, that is when both the neocortex.perc
and the mass is =0.

You could try to write down the model in the case you also have a
categorical variable

\subsection*{ANOVA vs Multiple linear
Regression}\label{anova-vs-multiple-linear-regression}
\addcontentsline{toc}{subsection}{ANOVA vs Multiple linear Regression}

\textbf{Student question (2021)}:

Hello I'm a bit confused when I have to test with an ANOVA and when with
a multiple linear regression. For example: The Yield of Hybrid Mais
example in lecture 6 compares the differences of means of the
groups(categorical variable). The Earthworm example is also a
categorical (explanatory) variable and there we test with a multiple
linear regression.

(we convert the categorical variable into a dummy variable but still
could I test there also with an ANOVA?)

\textbf{Answer given}:

Hi. ANOVA and (multiple) linear regression are both linear models that
investigate the relation between explanatory variables and a response
variable. In fact, ANOVA can be seen as a special case of linear
regression in which the \textbf{all} explanatory variables are
categorical.

In general, if there is a categorical variable (for instance, in a
ANOVA) we do not directly look at the \texttt{summary()} table (which
gives coefficient estimates and corresponding \emph{t}-test results) but
we look at the \texttt{anova()} table which gives the estimated SS, MS
and corresponding \emph{F}-test results. What is the difference? In the
\texttt{summary()} table each coefficient is separately tested with a
\emph{t}-test against the Null hypothesis that it is 0. For a
categorical variable with \emph{k} levels this means \emph{k-1} tests
(\emph{k-1} dummy variables) and if \(k\) is large we run into the
\href{https://en.wikipedia.org/wiki/Multiple_comparisons_problem}{multiple
comparison problem}. So, instead of testing each level of that
categorical variable separately we run 1 single test (the \emph{F}-test)
that tells us whether at least 2 levels are different from each other
(see slide 17 lecture 6). Afterwards we can look at the estimated
coefficients with \texttt{summary()}. For continuous variables we can
directly look at the \texttt{summary()} output. In the earthworm
example, there were tow explanatory variables, 1 categorical (Gattung)
and 1 continuous (Magenumf). For the former you first look at the
\texttt{anova()}table, for the latter you can directly look at the
\texttt{summary()} table.

\subsection*{Self test question 8 Spread and shape of
distributions}\label{self-test-question-8-spread-and-shape-of-distributions}
\addcontentsline{toc}{subsection}{Self test question 8 Spread and shape
of distributions}

\textbf{Student question (2021)}:

I find the question 8 of Spread and shape of distributions in the self
test not well formulated. The tail is the thinner part of the curve, so
if there is a long tail of high values it could mean that the data are
more concentrate in the lower part and therefore have a mena lower then
the median.

\textbf{Answer given}:

Hi! I can understand that it can be confusing. In statistics, what is
meant with ``long tail of high values'' (sometimes also called ``fat
tail'' or ``heavy tail'') is that there is a bigger probability of
getting (very) large values when compared to the reference distribution
(often a normal or an exponential distribution). See for instance
\href{https://en.wikipedia.org/wiki/Fat-tailed_distribution}{here}. So
the sentence ``A distribution of data with a long tail of high values''
always means that high values are more frequent than they would be in
the distribution of reference.

As you correctly identified, the mean is not a
\href{https://en.wikipedia.org/wiki/Robust_statistics\#Examples}{robust}
measure as its value is influenced by extreme values in the data, while
the median is more robust. So in the case of a distribution of data with
a long tail of high values, the mean is bigger than the median.

\subsection*{ggpair() graph exclusions}\label{ggpair-graph-exclusions}
\addcontentsline{toc}{subsection}{ggpair() graph exclusions}

\textbf{Student question (2021)}:

Dear BIO144 teammates. In unit 2, excercise 1, Question number 7, for
looking at relationships among variables we use the code ggpairs().
However, as you all have seen, it gives us all the graphs. It has some
cons including some of them are not needed for answering this question
and also the graphs are so small that it is hard to see the
relationship. Therefore I was wondering is there a way to exclude all
the graphs rather than the ones that we need?(the ones that show the
relationship of bodyfat with all other variables)

\textbf{Answer given}:

Hi! You could do it like this:

\begin{verbatim}
bodyfat\_dataset %>%
 select(bodyfat, age, abdomen, height, weight) %>% 
 ggpairs()
\end{verbatim}

\subsection*{Lecture4 page32}\label{lecture4-page32}
\addcontentsline{toc}{subsection}{Lecture4 page32}

\textbf{Student question (2021)}:

Could someone please look at our answers for the questions on this page?
we are not sure if our answers are correct and which other
interpretations are possible.

\begin{enumerate}
\def\labelenumi{\arabic{enumi}.}
\tightlist
\item
  x1 is an important explanatory variable because if only x1 is used,
  then R2 and the adjusted R2 are high and the p-value is small (what
  means that the slope for x1 is differs significantly from zero(?))
\item
  same answer for x2
\item
  In this model, we only need x1 or x2, not both. Reason: the R2 becomes
  not much higher if we use both compared to the situation where we only
  use one of them. x1 and x2 are also positively correlated: if x1
  increases, then x2 will increase too.
\item
  Interpretation: There is a positive correlation between y and each of
  the two x's: if x1 or x2 increases, then the catheter length will
  increase too.
\end{enumerate}

Thank you very much in advance for your answer.

\textbf{Answer given}:

\begin{itemize}
\tightlist
\item
  I think that your answer for questions 1 and 2 are fine (but see
  answer to question 3, which is relevant here as well).
\item
  \textbf{Question 3}: you will see this later in the course (so for now
  it is completely fine to just answer these questions as good as you
  can and if you want to you can ask, like you did), but the answer to
  this question is less straightforward and depends on our goals. In
  general, in regression there are 2 goals: \textbf{to predict} and
  \textbf{to explain}.
\item
  \textbf{To predict} means that we want to be able to predict the
  response variable y as good as we can, and we do not really care how
  we achieve this (i.e.~which variables we use) and we might just look
  at the adjusted r-squared value and pick the model with the highest
  value. There is a mistake on slide 31 and the adjusted r-squared for
  the model with both x1 and x2 is not shown, but it is 0.76. So in this
  case we would probably just take the model with just x2 (adjusted
  r-squared: 0.78)
\item
  \textbf{To explain} means that we are interested in the relation
  between the explanatory variables and the response variable (e.g.~what
  does a unit increase in x1 mean for y?). In this context, we can for
  instance fit two separate models for the two explanatory variables, or
  if we are only interested in one of the two, we just use that one. As
  I said, you will see this again later in the course.
\item
  \textbf{Question 4}: with this new information and what I already
  wrote above please try again on your own to interpret the model with
  both x1 and x2 (e.g.~are the slopes estimated to be significant?).
\end{itemize}

I hope this clears things up a little bit. You will see that things will
get clearer as you progress through the course as these things will come
up more than once. But feel free to ask for further clarifications.

\subsection*{Anova Degrees of Freedom}\label{anova-degrees-of-freedom}
\addcontentsline{toc}{subsection}{Anova Degrees of Freedom}

\textbf{Student question (2021)}:

Hello, I am quite confused about the degrees of freedom in Anova. In the
BC reading it says that the degree of freedom is the number of groups
one has. For the one-way anova example in the lecture, this would mean 4
so 20-4 = 16. However, in the slides it say the degree of freedom is
n-1, which would result in 19. When should I use which method or in
other words, when asked for the degrees of freedom of an Anova, which
number would be expected at an exam? Thanks for your help!

\textbf{Answer given}: There is more than one type of degree of freedom
involved, so first some theory: If you look at slide 18 lecture 6, you
can see that it is the total variability

SS\_\{total\} that has \textbf{n-1} degrees of freedom (i.e.~we need 1
degree of freedom to calculate it). In Anova, we partition this total
variability into the explained variability by the model

SS\_\{between\textasciitilde groups\} and into the residual variability

SS\_\{within\textasciitilde groups\} That is:

SS\_\{total\}=
SS\_\{between\textsubscript{groups\}+SS\_\{within}groups\} To calculate
the explained variability we need \textbf{g-1} degrees of freedom (with
\textbf{g} the number of groups) which leaves us with \textbf{n-1
-(g-1)=n-g} degrees of freedom for the residual variability. We then
test with the \emph{F}-test whether the explained variability is
significant. Under \emph{H0} the calculated test statistics

F=\frac{MS\_G}{MS\_E} follows a \emph{F}-distribution with degrees of
freedom \textbf{g-1} and \textbf{n-g}, i.e.~

F\sim F\_\{g-1,n-g\} so that's why we need those degrees of freedom.

Now, notice that the degrees of freedom used for the explained
variability is \textbf{g-1}, i.e.~one less than there are groups.
\textbf{In addition to this}, the intercept of the model is estimated as
well, which uses up another degree of freedom. Hence the total number of
degrees of freedom used by the anova model is \textbf{g-1+1=g}, and the
remaining degrees of freedom are \textbf{n-g} (or as in your example
\textbf{20-4=16}). Note that the intercept does not come into play in
the variance decomposition for the \emph{F}-test, hence it is not listed
in ANOVA table on slide 18.

I hope this clears things up a little.

\subsection*{Unit 10 Abalone Age 1}\label{unit-10-abalone-age-1}
\addcontentsline{toc}{subsection}{Unit 10 Abalone Age 1}

\textbf{Student question (2021)}:

I have a question concerning the interpretation of the data and I was a
bit confused with the terms of over-dispersion and under-dispersion. In
the interpretation of the data it says:

``So, interestingly, and rather unusually for biological data, the data
is \textbf{under-dispersed}, and''This is another example of a model
being \textbf{anti-conservative}.''

And in the summary table we can see quite \textbf{low p-values}.

But in the lecture slides (lecture 10, slide 38) we learned that ``When
there is unaccounted over-dispersion, the p-values that are calculated
are usually too small!'' And on slide 41, which talks about
\textbf{under-dispersion}, it says ``In that case, your p-values are
usually \textbf{too large}, that is, the results are
\textbf{conservative}''

So I don't understand the correct connection between
over/under-dispersion, low/high p-values and (anti-)/conservative.

\textbf{Answer given}:

It is just a matter of slightly unlucky placement: the sentence ``This
is another example of a model being anti-conservative.'' comes after a
horizontal line (denoting a new topic/paragraph/etc.) and is preceded by
the sentence ``\textbf{So the poisson glm had fewer significant terms
than the lm.}'' (\textbf{Note}: this refers to the old version on
openedx, and not to the version on Olat). Hence it refers to that
sentence and has nothing to do with whether the model is over- or
underdispersed. In fact, there are many reasons for why something
(p-value, confidence interval, etc.) can be (anti-)conservative, one
being that a wrong model is used: the lm is the wrong model (because
it's count data) and it produced smaller p-value than the glm, thus
\emph{in this case} (!) the lm is anti-conservative. What is written
about the dispersion parameter in the slides is correct.

\subsection*{Unit 10 Abalone Age 2}\label{unit-10-abalone-age-2}
\addcontentsline{toc}{subsection}{Unit 10 Abalone Age 2}

\textbf{Student question (2021)}:

I have a question concerning the interpretation of the data and I was a
bit confused with the terms of over-dispersion and under-dispersion. In
the interpretation of the data it says:

``So, interestingly, and rather unusually for biological data, the data
is \textbf{under-dispersed}, and''This is another example of a model
being \textbf{anti-conservative}.''

And in the summary table we can see quite \textbf{low p-values}.

But in the lecture slides (lecture 10, slide 38) we learned that ``When
there is unaccounted over-dispersion, the p-values that are calculated
are usually too small!'' And on slide 41, which talks about
\textbf{under-dispersion}, it says ``In that case, your p-values are
usually \textbf{too large}, that is, the results are
\textbf{conservative}''

So I don't understand the correct connection between
over/under-dispersion, low/high p-values and (anti-)/conservative.

\textbf{Answer given}:

It is just a matter of slightly unlucky placement: the sentence ``This
is another example of a model being anti-conservative.'' comes after a
horizontal line (denoting a new topic/paragraph/etc.) and is preceded by
the sentence ``\textbf{So the poisson glm had fewer significant terms
than the lm.}'' (\textbf{Note}: this refers to the old version on
openedx, and not to the version on Olat). Hence it refers to that
sentence and has nothing to do with whether the model is over- or
underdispersed. In fact, there are many reasons for why something
(p-value, confidence interval, etc.) can be (anti-)conservative, one
being that a wrong model is used: the lm is the wrong model (because
it's count data) and it produced smaller p-value than the glm, thus
\emph{in this case} (!) the lm is anti-conservative. What is written
about the dispersion parameter in the slides is correct.

\subsection*{Chr vs Factors}\label{chr-vs-factors}
\addcontentsline{toc}{subsection}{Chr vs Factors}

\textbf{Student question (2021)}:

When I import my data, I often have ``chr'' rather than ``factor''. This
was the case for country and continent in the healthcare\_financing.csv.

Should I do something specific to get the data directly as factor? If
not, is there a way to change all them at once, rather than to change
them one after the other using as.factor?

\textbf{Answer given}:

As of R 4.0.0. variables are read into R as characters instead of
factors. Ahead of this course it was decided that we keep it like this
because linear models can be fitted with both types. So it is suggested
that you do not change the variables to factors (in fact, the only time
you might/will want to change them is when you want to change the levels
within a factor\ldots{} and for now at least this is not needed!). If
nevertheless you want to change characters to factors, it might actually
be good to do that one variable at a time as you will only convert only
want you need to convert and it gives you more control over it.

PS: if you really want to convert all characters to factors, you can do
this (but again, no need to do it!)

\begin{verbatim}
dd <- dd %>% mutate(across(where(is\_character),as\_factor)) 
\end{verbatim}

\subsection*{Changing reference in earthworm video
example}\label{changing-reference-in-earthworm-video-example}
\addcontentsline{toc}{subsection}{Changing reference in earthworm video
example}

\textbf{Student question (2021)}:

In the earthworm analysis of the correlation between log.gewicht and
Gattung. Is it correct to think that if the reference Gattung had been
``N'' , which seems to have a similar mean than Oc, then the p-value in
the linear regression model would only have been significant for L but
would not have been significant for Oc? I am assuming that the means for
N and Oc are statistically similar.

\textbf{Answer given}:

Yes, exactly. You can change the reference level with the following
code, then run the model with that new Gattung variable, and check the
summary table. (You must install the forcats package to use the
fct\_relevel function.)

library(forcats)

\begin{verbatim}
 dd <- dd %>%
     mutate(Gattung\_refN\_ = fct\_relevel(Gattung, "N", after = 0))
\end{verbatim}

\subsection*{Decomposing R\^{}2}\label{decomposing-r2}
\addcontentsline{toc}{subsection}{Decomposing R\^{}2}

Student question (2021):

I have a question concerning the decomposition of R\^{}2. In the lecture
we learned that we should calculate the relative importance with the
package relaimpo and the function calc.relimp. However, in the IC
material, we learned how to calculate it more manually (R\^{}2 of
model\_both - R\^{}2 of model\_weight). These two methods do not produce
the same result so when should we use which method? Or did I understand
something incorrectly in general?

Answer given:

There is more than one way to calculate relative importances of
variables and the various methods differ in the the produced results. If
you will be asked to calculate them, you will be told which approach to
use.

\bookmarksetup{startatroot}

\chapter*{Discussion Forum}\label{discussion-forum}
\addcontentsline{toc}{chapter}{Discussion Forum}

\markboth{Discussion Forum}{Discussion Forum}

Please use the discussion forum in OLAT, here.

\bookmarksetup{startatroot}

\chapter*{Unit 1}\label{unit-1}
\addcontentsline{toc}{chapter}{Unit 1}

\markboth{Unit 1}{Unit 1}

\section*{Introduction}\label{introduction-1}
\addcontentsline{toc}{section}{Introduction}

\markright{Introduction}

Welcome to Week 1 of BIO144! Week 1 is a bit different from other weeks
because its focuses on introducing how the course is going to work,
collecting some data for us to use, and reviewing some material that it
would be good for you to remember from previous courses (or learn here
if you don't know it yet!)

In subsequent weeks we will focus more on the data analysis concepts and
practice.

However, just like all weeks, there is:

\begin{itemize}
\tightlist
\item
  A chapter of the course book.
\item
  A ``lecture'' on Monday.
\item
  Some homework to ideally do between Monday and the practical on
  Thursday or Friday.
\item
  A practical session on Thursday or Friday.
\item
  A weekly quiz to check your understanding.
\end{itemize}

\section*{Course book}\label{course-book}
\addcontentsline{toc}{section}{Course book}

\markright{Course book}

This week we will be covering Chapter 1 of
\href{https://opetchey.github.io/BIO144_Course_Book/}{the course book}.

\section*{Lecture}\label{lecture}
\addcontentsline{toc}{section}{Lecture}

\markright{Lecture}

In the lecture you'll be asked to record some reaction time data about
yourself. Use this google form to record the data: Link to Reaction Time
Data Form

Go to this web page for a tool to record your reaction times: Link to
Reaction Time Tool

The resulting dataset can be access and or downloaded from this web
page: Link to Reaction Time Dataset

\section*{Homework, Practical, and Weekly
Quiz}\label{homework-practical-and-weekly-quiz}
\addcontentsline{toc}{section}{Homework, Practical, and Weekly Quiz}

\markright{Homework, Practical, and Weekly Quiz}

These are on this web page: Link to Homework, Practical, and Weekly Quiz

\bookmarksetup{startatroot}

\chapter*{Unit 2}\label{unit-2}
\addcontentsline{toc}{chapter}{Unit 2}

\markboth{Unit 2}{Unit 2}

\section*{Introduction}\label{introduction-2}
\addcontentsline{toc}{section}{Introduction}

\markright{Introduction}

Welcome to Week 2 of BIO144!

\section*{Course book}\label{course-book-1}
\addcontentsline{toc}{section}{Course book}

\markright{Course book}

This week we will be covering Chapter 2 of
\href{https://opetchey.github.io/BIO144_Course_Book/}{the course book}.

\section*{Lecture}\label{lecture-1}
\addcontentsline{toc}{section}{Lecture}

\markright{Lecture}

Please attend the lecture. If for any reason you can't, please watch the
podcast, which is available on OLAT.

\section*{Homework, Practical, and Weekly
Quiz}\label{homework-practical-and-weekly-quiz-1}
\addcontentsline{toc}{section}{Homework, Practical, and Weekly Quiz}

\markright{Homework, Practical, and Weekly Quiz}

These are on this web page: Link to Homework, Practical, and Weekly Quiz

\bookmarksetup{startatroot}

\chapter*{Unit 3}\label{unit-3}
\addcontentsline{toc}{chapter}{Unit 3}

\markboth{Unit 3}{Unit 3}

\section*{Introduction}\label{introduction-3}
\addcontentsline{toc}{section}{Introduction}

\markright{Introduction}

Welcome to Week 3 of BIO144!

\section*{Course book}\label{course-book-2}
\addcontentsline{toc}{section}{Course book}

\markright{Course book}

This week we will be covering Chapter 3 of
\href{https://opetchey.github.io/BIO144_Course_Book/}{the course book}.

\section*{Lecture}\label{lecture-2}
\addcontentsline{toc}{section}{Lecture}

\markright{Lecture}

\section*{Homework, Practical, and Weekly
Quiz}\label{homework-practical-and-weekly-quiz-2}
\addcontentsline{toc}{section}{Homework, Practical, and Weekly Quiz}

\markright{Homework, Practical, and Weekly Quiz}

These are on this web page: Link to Homework, Practical, and Weekly Quiz

\bookmarksetup{startatroot}

\chapter*{Unit 4}\label{unit-4}
\addcontentsline{toc}{chapter}{Unit 4}

\markboth{Unit 4}{Unit 4}

\section*{Introduction}\label{introduction-4}
\addcontentsline{toc}{section}{Introduction}

\markright{Introduction}

Welcome to Week 4 of BIO144!

\section*{Course book}\label{course-book-3}
\addcontentsline{toc}{section}{Course book}

\markright{Course book}

This week we will be covering Chapter 4 of
\href{https://opetchey.github.io/BIO144_Course_Book/}{the course book}.

\section*{Lecture}\label{lecture-3}
\addcontentsline{toc}{section}{Lecture}

\markright{Lecture}

\section*{Homework, Practical, and Weekly
Quiz}\label{homework-practical-and-weekly-quiz-3}
\addcontentsline{toc}{section}{Homework, Practical, and Weekly Quiz}

\markright{Homework, Practical, and Weekly Quiz}

These are on this web page: Link to Homework, Practical, and Weekly Quiz

\bookmarksetup{startatroot}

\chapter*{Unit 5}\label{unit-5}
\addcontentsline{toc}{chapter}{Unit 5}

\markboth{Unit 5}{Unit 5}

\section*{Introduction}\label{introduction-5}
\addcontentsline{toc}{section}{Introduction}

\markright{Introduction}

Welcome to Week 5 of BIO144!

\section*{Course book}\label{course-book-4}
\addcontentsline{toc}{section}{Course book}

\markright{Course book}

This week we will be covering Chapter 5 of
\href{https://opetchey.github.io/BIO144_Course_Book/}{the course book}.

\section*{Lecture}\label{lecture-4}
\addcontentsline{toc}{section}{Lecture}

\markright{Lecture}

\section*{Homework, Practical, and Weekly
Quiz}\label{homework-practical-and-weekly-quiz-4}
\addcontentsline{toc}{section}{Homework, Practical, and Weekly Quiz}

\markright{Homework, Practical, and Weekly Quiz}

These are on this web page: Link to Homework, Practical, and Weekly Quiz

\bookmarksetup{startatroot}

\chapter*{Unit 6}\label{unit-6}
\addcontentsline{toc}{chapter}{Unit 6}

\markboth{Unit 6}{Unit 6}

\section*{Introduction}\label{introduction-6}
\addcontentsline{toc}{section}{Introduction}

\markright{Introduction}

Welcome to Week 6 of BIO144!

\section*{Course book}\label{course-book-5}
\addcontentsline{toc}{section}{Course book}

\markright{Course book}

This week we will be covering Chapter 6 of
\href{https://opetchey.github.io/BIO144_Course_Book/}{the course book}.

\section*{Lecture}\label{lecture-5}
\addcontentsline{toc}{section}{Lecture}

\markright{Lecture}

\section*{Homework, Practical, and Weekly
Quiz}\label{homework-practical-and-weekly-quiz-5}
\addcontentsline{toc}{section}{Homework, Practical, and Weekly Quiz}

\markright{Homework, Practical, and Weekly Quiz}

These are on this web page: Link to Homework, Practical, and Weekly Quiz

\bookmarksetup{startatroot}

\chapter*{Unit 7}\label{unit-7}
\addcontentsline{toc}{chapter}{Unit 7}

\markboth{Unit 7}{Unit 7}

\section*{Introduction}\label{introduction-7}
\addcontentsline{toc}{section}{Introduction}

\markright{Introduction}

Welcome to Week 7 of BIO144!

\section*{Course book}\label{course-book-6}
\addcontentsline{toc}{section}{Course book}

\markright{Course book}

This week we will be covering Chapter 7 of
\href{https://opetchey.github.io/BIO144_Course_Book/}{the course book}.

\section*{Lecture}\label{lecture-6}
\addcontentsline{toc}{section}{Lecture}

\markright{Lecture}

\section*{Homework, Practical, and Weekly
Quiz}\label{homework-practical-and-weekly-quiz-6}
\addcontentsline{toc}{section}{Homework, Practical, and Weekly Quiz}

\markright{Homework, Practical, and Weekly Quiz}

These are on this web page: Link to Homework, Practical, and Weekly Quiz

\bookmarksetup{startatroot}

\chapter*{Unit 8}\label{unit-8}
\addcontentsline{toc}{chapter}{Unit 8}

\markboth{Unit 8}{Unit 8}

\section*{Introduction}\label{introduction-8}
\addcontentsline{toc}{section}{Introduction}

\markright{Introduction}

Welcome to Week 8 of BIO144!

\section*{Course book}\label{course-book-7}
\addcontentsline{toc}{section}{Course book}

\markright{Course book}

This week we will be covering Chapter 8 of
\href{https://opetchey.github.io/BIO144_Course_Book/}{the course book}.

\section*{Lecture}\label{lecture-7}
\addcontentsline{toc}{section}{Lecture}

\markright{Lecture}

\section*{Homework, Practical, and Weekly
Quiz}\label{homework-practical-and-weekly-quiz-7}
\addcontentsline{toc}{section}{Homework, Practical, and Weekly Quiz}

\markright{Homework, Practical, and Weekly Quiz}

These are on this web page: Link to Homework, Practical, and Weekly Quiz

\bookmarksetup{startatroot}

\chapter*{Unit 9}\label{unit-9}
\addcontentsline{toc}{chapter}{Unit 9}

\markboth{Unit 9}{Unit 9}

\section*{Introduction}\label{introduction-9}
\addcontentsline{toc}{section}{Introduction}

\markright{Introduction}

Welcome to Week 9 of BIO144!

\section*{Course book}\label{course-book-8}
\addcontentsline{toc}{section}{Course book}

\markright{Course book}

This week we will be covering Chapter 9 of
\href{https://opetchey.github.io/BIO144_Course_Book/}{the course book}.

\section*{Lecture}\label{lecture-8}
\addcontentsline{toc}{section}{Lecture}

\markright{Lecture}

\section*{Homework, Practical, and Weekly
Quiz}\label{homework-practical-and-weekly-quiz-8}
\addcontentsline{toc}{section}{Homework, Practical, and Weekly Quiz}

\markright{Homework, Practical, and Weekly Quiz}

These are on this web page: Link to Homework, Practical, and Weekly Quiz

\bookmarksetup{startatroot}

\chapter*{Unit 10}\label{unit-10}
\addcontentsline{toc}{chapter}{Unit 10}

\markboth{Unit 10}{Unit 10}

\section*{Introduction}\label{introduction-10}
\addcontentsline{toc}{section}{Introduction}

\markright{Introduction}

Welcome to Week 10 of BIO144!

\section*{Course book}\label{course-book-9}
\addcontentsline{toc}{section}{Course book}

\markright{Course book}

This week we will be covering Chapter 10 of
\href{https://opetchey.github.io/BIO144_Course_Book/}{the course book}.

\section*{Lecture}\label{lecture-9}
\addcontentsline{toc}{section}{Lecture}

\markright{Lecture}

\section*{Homework, Practical, and Weekly
Quiz}\label{homework-practical-and-weekly-quiz-9}
\addcontentsline{toc}{section}{Homework, Practical, and Weekly Quiz}

\markright{Homework, Practical, and Weekly Quiz}

These are on this web page: Link to Homework, Practical, and Weekly Quiz

\bookmarksetup{startatroot}

\chapter*{Unit 11}\label{unit-11}
\addcontentsline{toc}{chapter}{Unit 11}

\markboth{Unit 11}{Unit 11}

\section*{Introduction}\label{introduction-11}
\addcontentsline{toc}{section}{Introduction}

\markright{Introduction}

Welcome to Week 11 of BIO144!

\section*{Course book}\label{course-book-10}
\addcontentsline{toc}{section}{Course book}

\markright{Course book}

This week we will be covering Chapter 11 of
\href{https://opetchey.github.io/BIO144_Course_Book/}{the course book}.

\section*{Lecture}\label{lecture-10}
\addcontentsline{toc}{section}{Lecture}

\markright{Lecture}

\section*{Homework, Practical, and Weekly
Quiz}\label{homework-practical-and-weekly-quiz-10}
\addcontentsline{toc}{section}{Homework, Practical, and Weekly Quiz}

\markright{Homework, Practical, and Weekly Quiz}

These are on this web page: Link to Homework, Practical, and Weekly Quiz

\bookmarksetup{startatroot}

\chapter*{Unit 12}\label{unit-12}
\addcontentsline{toc}{chapter}{Unit 12}

\markboth{Unit 12}{Unit 12}

\section*{Introduction}\label{introduction-12}
\addcontentsline{toc}{section}{Introduction}

\markright{Introduction}

Welcome to Week 12 of BIO144!

\section*{Course book}\label{course-book-11}
\addcontentsline{toc}{section}{Course book}

\markright{Course book}

This week we will be covering Chapter 12 of
\href{https://opetchey.github.io/BIO144_Course_Book/}{the course book}.

\section*{Lecture}\label{lecture-11}
\addcontentsline{toc}{section}{Lecture}

\markright{Lecture}

\section*{Homework, Practical, and Weekly
Quiz}\label{homework-practical-and-weekly-quiz-11}
\addcontentsline{toc}{section}{Homework, Practical, and Weekly Quiz}

\markright{Homework, Practical, and Weekly Quiz}

These are on this web page: Link to Homework, Practical, and Weekly Quiz


\backmatter


\end{document}
